% 底本源行 1--191。
\markboth{第 XVI 讲. 幺幂秩为零的群}{第 XVI 讲. 幺幂秩为零的群}
\setcounter{footnote}{0}
\SGARefTarget{sga3:xvi:chapter:xvi}\chapter*{第 XVI 讲\\幺幂秩为零的群}
\addcontentsline{toc}{chapter}{第 XVI 讲. 幺幂秩为零的群}
\markboth{第 XVI 讲. 幺幂秩为零的群}{第 XVI 讲. 幺幂秩为零的群}

\begin{center}
  \textit{M. Raynaud}\textsuperscript{(*)}
\end{center}
\begingroup
\renewcommand{\thefootnote}{0}
\footnotetext{2008 年 12 月 1 日的 xy 版本。}
\renewcommand{\thefootnote}{*}
\footnotetext{参见第 XV 讲第 1 页的注。}
\endgroup

\SGARefTarget{sga3:xvi:section:1}\section*{1. 一个浸入判据}
\addcontentsline{toc}{section}{1. 一个浸入判据}

\SGARefTarget{sga3:xvi:subsection:1:1}%
\subsection*{1.1. 不是浸入的群预概形单态射的例子}
\addcontentsline{toc}{subsection}{\SGARefLink{sga3:xvi:subsection:1:1}{1.1}. 不是浸入的群预概形单态射的例子}

我们将构造预概形 \(S\)、两个 \(S\)-群预概形 \(G\)、\(H\)，以及不是浸入的 \(S\)-群单态射
\[
  u:G\longrightarrow H.
\]
群 \(G\)、\(H\) 在 \(S\) 上有限表示，在 \(S\) 上平坦，且 \(G\) 甚至在 \(S\) 上光滑（参见第 VIII 讲 \SGARefLink{sga3:viii:lemma:7:1}{7.1} 前的一段及注 \((*)\)；亦见 XVII，附录 III，4）。

先取 \(S\) 为混合特征、剩余特征为 \(2\) 的离散赋值环 \(A\) 的谱。设 \(t\) 为 \(S\) 的一般点，\(s\) 为闭点。

\medskip
\noindent\textit{a)} 从常群
\[
  G'=(\mathbb Z/2\mathbb Z)_S
\]
中去掉闭纤维 \((\mathbb Z/2\mathbb Z)_s\) 中不同于单位元的点，取所得开子群为 \(G\)；取 \(H\) 为二次单位根的乘法型群 \((\mu_2)_S\)（\SGARefLink{sga3:i:paragraph:4-4-4}{第 I 讲，4.4.4}）。由 \(S\) 的选取方式，\(H_t\) 同构于 \((\mathbb Z/2\mathbb Z)_t\)，而 \(H_s\) 为纯不可分群。由 \(H\) 的截面 \((-1)\) 定义显然的态射 \(G'\to H\)，从而给出态射 \(u:G\to H\)，它为单态射而不是浸入（否则必为同构；参见\SGARefLink{sga3:viii:section:7}{第 VIII 讲，7}）。此处\SGARefLink{sga3:viii:proposition:7:9}{第 VIII 讲 7.9} 不适用，因为 \({}_{2}G=G\) 在 \(S\) 上不有限。

\medskip
\noindent\textit{b)} 设 \(K\) 为将单位群的闭纤维的唯一点“加倍”而得到的平展、非分离 \(S\)-群预概形。取 \(H\) 为乘积群
\[
  (\mu_2)_S\times_S K.
\]
设 \(a\)（分别地，\(b\)）为 \(\mu_2\)（分别地，\(K\)）在 \(S\) 上唯一的不同于单位截面的截面，设 \(h=(a,b)\) 为 \(H\) 的相应截面。由 \(h\) 定义一个 \(S\)-群态射
\[
  u:G=(\mathbb Z/2\mathbb Z)_S\longrightarrow H.
\]
显然 \(u\) 为单态射而不是浸入。此处\SGARefLink{sga3:viii:proposition:7:9}{第 VIII 讲 7.9} 不适用，因为 \({}_{2}H=H\) 在 \(S\) 上不分离。

\medskip
\noindent\SGAInlineRefTarget{sga3:xvi:subsection:1:1:item:c}\textit{c)} 下面的例子更有意思，其中 \(G\) 是纤维连通的光滑群（甚至 \(G=(\Ga)_S\)）；参见\SGARefLink{sga3:viii:remark:7:10}{第 VIII 讲，7.10}。

取 \(S\) 为等特征 \(2\) 的离散赋值环 \(A\) 的谱，设 \(\pi\) 为 \(A\) 的一个素元，\(\mathfrak m\) 为 \(A\) 的极大理想。

考虑加法群 \((\Ga)_S\)、环为 \(A[X,Y]\) 的乘积群 \((\Ga\times\Ga)_S\)，并用下式定义闭子群 \(G_1\)：
\begin{equation}
  X^2+X-\pi Y=0.
  \tag{1}\label{sga3:xvi:eq:1}
\end{equation}
群 \(G_1\) 在 \(S\) 上光滑；其一般纤维同构于 \((\Ga)_t\)，闭纤维同构于 \((\Ga)_s\) 与 \((\mathbb Z/2\mathbb Z)_s\) 的乘积。对 \((\mathbb Z/2\mathbb Z)_s\) 因子，可取 \((G_1)_s\) 中由下式定义的子群 \(N_s\)：
\begin{equation}
  Y=X^2+X=0.
  \tag{2}\label{sga3:xvi:eq:2}
\end{equation}

设 \(N'\)、\(N''\) 为 \(G_1\) 的两个同构于 \((\mathbb Z/2\mathbb Z)_S\) 的子群概形，其闭纤维同为 \(N_s\)，而一般纤维不同。因此 \(N'\)、\(N''\) 由 \(G_1\) 在 \(S\) 上的两个截面定义，其坐标分别为 \((a',b')\)、\((a'',b'')\)，使
\begin{equation}
  \begin{cases}
    a'^2+a'-\pi b'=a''^2+a''-\pi b''=0,\\
    a'\text{ 与 }a''\equiv 1\pmod{\mathfrak m},
      \qquad a'\ne a'',\\
    b'\text{ 与 }b''\equiv 0\pmod{\mathfrak m}.
  \end{cases}
  \tag{3}\label{sga3:xvi:eq:3}
\end{equation}
（例如可取 \((1,0)\)、\((1+\pi^2,\pi+\pi^3)\)。）

商群
\[
  G'=G_1/N'
  \qquad\text{与}\qquad
  G''=G_1/N''
\]
可表（\SGARefLink{thm:sga3-V-4.1}{第 V 讲，4.1}）。设 \(u\) 为典范态射
\[
  u:G_1\longrightarrow G'\times_S G'',
\]
设 \(H\) 为 \(G'\times_S G''\) 中等于 \(u_t((G_1)_t)\) 在 \(G'\times_S G''\) 中的概形闭包的子群概形，于是 \(u\) 通过 \(H\) 分解。\(u\) 的核为
\[
  K=N'\cap N'',
\]
其一般纤维为单位群，闭纤维等于 \(N_s\)；特别地，\(K\) 在 \(S\) 上不平坦。因此在一般纤维上，\(u_t\) 将 \((G_1)_t\) 同构到 \(H_t\)。另一方面，\(H_s\) 不光滑。事实上，若 \(H_s\) 光滑，则由于 \(u_s\) 为核有限的态射，且 \((G_1)_s\)、\(H_s\) 的维数相同，均为 \(1\)，\(u_s\) 就平坦。由于 \(G_1\)、\(H\) 平坦，\(u\) 此时也平坦，因而 \(\Ker u\) 在 \(S\) 上平坦，矛盾。显然，\(u\) 在 \(G_1\) 的单位连通分支 \(G\) 上的限制为单态射：\(K\cap G\) 的纤维为单位群，故 \(K\cap G\) 为单位群（第 VI\textsubscript{B} 讲，\SGARefLink{lem:sga3-VIB-2.9}{2.9}）。但它不是浸入，否则它为开浸入，从而 \(H_s\) 光滑。此处\SGARefLink{sga3:viii:proposition:7:9}{第 VIII 讲 7.9} 不适用，因为 \({}_{2}G=G\) 在 \(S\) 上不有限。

对于喜欢方程的读者，在上例中，可取 \(H\) 为 \((\Ga\times_S\Ga)_S\) 中环为 \(A[V,W]/(F)\) 的闭子群，其中
\[
  F=(a''b'-b''a')(a''^2V-a'^2W)-(a''V-a'W)^2
\]
（其中 \(a',a'',b',b''\) 满足 \eqref{sga3:xvi:eq:3}）。对 \(G\)，取环为 \(A[T]\) 的群 \((\Ga)_S\)；对态射 \(u:G\to H\)，取由下式定义的 \(S\)-态射：
\[
  V\longmapsto a'(a'T+\pi T^2),
  \qquad
  W\longmapsto a''(a''T+\pi T^2).
\]

\medskip
\noindent\SGARefTarget{sga3:xvi:remark:1:2}%
\textbf{注 1.2。}---上述构造受到 Koizumi--Shimura 方法的启发。\footnote{编者注：S.~Koizumi 和 G.~Shimura，``Specialization of abelian varieties,'' \emph{Sci.\ Papers Coll.\ Gen.\ Ed.\ Univ.\ Tokyo} 9（1959），187--211。}当 \(S\) 为混合特征环的谱时，此方法不适用，因为此时 \(G\) 的有限阶点“过于刚性”。

\SGARefTarget{sga3:xvi:subsection:1:3}%
\subsection*{1.3. 浸入判据的陈述}
\addcontentsline{toc}{subsection}{1.3. 浸入判据的陈述}

\medskip
\noindent\SGARefTarget{sga3:xvi:theorem:1:3}%
\textbf{定理 1.3。}---\textit{设 \(S\) 为局部 noether 预概形，\(G\) 为 \(S\) 上光滑且有限型的 \(S\)-群预概形，并在 fpqc 拓扑下局部具有 Cartan 子群（\SGARefLink{sga3:xv:definition:6-1}{第 XV 讲，6.1} 及 \SGARefLink{sga3:xv:theorem:7-3}{7.3}~(i)）；设 \(H\) 为局部有限型的 \(S\)-群预概形，\(u:G\to H\) 为 \(S\)-群单态射。则：}
\begin{enumerate}
\item[\textit{a)}] \textit{若 \(G\) 的纤维连通，则 \(u\) 为浸入，当且仅当对每个形如剩余域代数闭的完备离散赋值环之谱的 \(S\)-概形 \(S'\)，以及 \(G_{S'}\) 的每个 Cartan 子群 \(C\)，\(u_{S'}\) 在 \(C\) 上的限制均为浸入。}
\item[\textit{b)}] \textit{态射 \(u\) 为浸入（分别为闭浸入），当且仅当对每个上述 \(S'\) 及 \(G_{S'}\) 的单位分支 \((G_{S'})^0\) 的每个 Cartan 子群 \(C\)，\(u_{S'}\) 在 \(C\) 于 \(G_{S'}\) 中的正规化子 \(N\)（\SGARefLink{sga3:xi:corollary:6:11}{第 XI 讲，6.11}）上的限制为浸入（分别为闭浸入）。}
\end{enumerate}

在证明 \SGARefLink{sga3:xvi:theorem:1:3}{1.3} 之前，先陈述几个应用。

\medskip
\noindent\SGARefTarget{sga3:xvi:corollary:1:4}%
\textbf{推论 1.4。}---\textit{设 \(S\) 为预概形，\(G\) 为 \(S\) 上光滑且有限表示的 \(S\)-群预概形，幺幂秩为零（\SGARefLink{sga3:xv:definition:6:1:ter}{第 XV 讲，6.1 ter}），且在 fpqc 拓扑下局部具有极大环面；设 \(H\) 为 \(S\)-群预概形，\(u:G\to H\) 为 \(S\)-群单态射。再假设 \(H\) 在 \(S\) 上有限表示，或 \(S\) 局部 noether 且 \(H\) 局部有限型。则：}
\begin{enumerate}
\item[\textit{a)}] \textit{若 \(G\) 的纤维连通，则 \(u\) 为浸入。}
\item[\textit{b)}] \textit{若 \(H\) 在 \(S\) 上分离，且对每个 \(S\) 上的 \(S'\) 及 \(G_{S'}\) 的每个极大环面 \(T\)，Weyl 群}
\[
  W=
  \Norm_{G_{S'}}(T)\big/\Centr_{G_{S'}^0}(T)
\]
\textit{可表（当 \(G\) 的纤维连通时，此条件总成立；\SGARefLink{sga3:xv:proposition:7-1}{第 XV 讲，7.1}~(iv)），且在 \(S\) 上有限，则 \(u\) 为闭浸入。}
\end{enumerate}

\medskip
\noindent\SGARefTarget{sga3:xvi:corollary:1:5}%
\textbf{推论 1.5。}---\textit{设 \(S\) 为预概形，\(G\) 为 \(S\) 上光滑且有限表示、纤维连通的 \(S\)-群预概形，\(H\) 为 \(S\)-群预概形，\(u:G\to H\) 为 \(S\)-群单态射。假设 \(H\) 在 \(S\) 上有限表示，或 \(S\) 局部 noether 且 \(H\) 局部有限型。则：}
\begin{enumerate}
\item[\textit{a)}] \textit{若 \(G\) 约化，即在 \(S\) 上仿射且纤维约化（参见第 XIX 讲），则 \(u\) 为浸入；若 \(H\) 分离，则为闭浸入。}
\item[\textit{b)}] \textit{若 \(G\) 的纤维仿射、可解、连通且幺幂秩为零，则 \(u\) 为浸入；若 \(H\) 分离，则为闭浸入。}
\end{enumerate}

\medskip
\noindent\textit{由 \SGARefLink{sga3:xvi:theorem:1:3}{1.3} 证明 \SGARefLink{sga3:xvi:corollary:1:4}{1.4}。}

\smallskip
\noindent\textit{化归到 noether 情形。}
若 \(S\) 局部 noether，由于待证性质在 \(S\) 上是局部的，化归立即成立。在第二种情形下，\(G\) 和 \(H\) 在 \(S\) 上有限表示。缩小 \(S\) 后，可设 \(S\) 仿射。由第 VI\textsubscript{B} 讲第 10 节，存在 noether 概形 \(S_0\)、在 \(S_0\) 上光滑（在情形 a) 中纤维连通）且在 \(S_0\) 上有限型的 \(S_0\)-群预概形 \(G_0\)、有限型的 \(S_0\)-群预概形 \(H_0\)（在情形 b) 中分离），以及 \(S_0\)-群态射
\[
  u_0:G_0\longrightarrow H_0
\]
使得 \(G,H,u\) 由 \(G_0,H_0,u_0\) 经基扩张 \(S\to S_0\) 得到。\(u\) 是单态射，意即 \(\Ker u\) 为单位群；因此可设 \(u_0\) 为单态射。

由于 \(G\) 的幺幂秩 \(\rho_u=0\)，且在 fpqc 拓扑下局部具有极大环面，\(G\) 的纤维的 Abel 秩 \(\rho_{ab}\) 局部常值（\SGARefLink{sga3:xv:corollary:8-18}{第 XV 讲，8.18}）。而 \(\rho_u\) 和 \(\rho_{ab}\) 是局部可构造函数（第 XV 讲，6.3 bis）。标准论证（参见 EGA IV，8.3.4）表明，可以选取 \(S_0\) 和 \(G_0\)，使 \(G_0\) 的纤维的幺幂秩（分别为 Abel 秩）为零（分别为局部常值）。于是 \(G_0\) 在 fpqc 拓扑下局部具有极大环面（\SGARefLink{sga3:xv:corollary:8-18}{第 XV 讲，8.18}），且 \(G_0\) 的极大环面函子 \(\mathcal{TM}_{G_0}\) 由 \(S_0\) 上有限型的 \(S_0\)-概形 \(X_0\) 表示（\SGARefLink{sga3:xv:theorem:8-15}{第 XV 讲，8.15}）。设 \(T_0\) 为 \((G_0)_{X_0}\) 的“泛”极大环面，令 \(X=X_0\times_{S_0}S\)，并令 \(T=T_0\times_{S_0}S\)，即 \(G\) 的泛极大环面。由假设，在情形 b) 中，相对于 \(T\) 的 Weyl 群可表，且在 \(X\) 上有限。这两个性质与预概形的投射极限相容（第 VI\textsubscript{B} 讲，10.1~iii 及 EGA IV，8.10.5）。因此可以选取 \(S_0\)，使 \(T_0\) 的 Weyl 群在 \(X_0\) 上有限。在这些条件下，为证明 \SGARefLink{sga3:xvi:corollary:1:4}{1.4}，可将 \(S,G,u,H\) 换为 \(S_0,G_0,u_0,H_0\)，从而可设 \(S\) 为 noether 概形。

使用 \SGARefLink{sga3:xvi:theorem:1:3}{1.3} 给出的赋值判据，化归到 \(S\) 为离散赋值环的谱且 \(G^0\) 具有 Cartan 子群 \(C\) 的情形。

\smallskip
\noindent\textit{\SGARefLink{sga3:xvi:corollary:1:4}{1.4} a) 的证明。}
必须证明 \(u\) 在 \(C\) 上的限制为浸入（\SGARefLink{sga3:xvi:theorem:1:3}{1.3}~a)），这将问题化归到 \(G=C\) 的纤维连通且幂零的情形。标准化归（参见\SGARefLink{sga3:viii:lemma:7:1}{第 VIII 讲，7.1} 的证明）进一步化归到 \(H\) 在 \(S\) 上平坦、\(u\) 在泛纤维上为同构、\(u\) 在闭纤维上为底空间的同构的情形。此时群 \(H\) 的纤维连通，因而在 \(S\) 上分离（第 VI\textsubscript{B} 讲，\SGARefLink{prop:sga3-VIB-5.2}{5.2}）。暂且承认下述引理。

\medskip
\noindent\SGARefTarget{sga3:xvi:lemma:1:6}%
\textbf{引理 1.6。}---\textit{设 \(S\) 为预概形，\(G\) 为在 \(S\) 上光滑、纤维连通且幂零、幺幂秩为零的 \(S\)-群预概形。则：}
\begin{enumerate}
\item[\textit{i)}] \textit{\(G\) 交换。}
\item[\textit{ii)}] \textit{对每个整数 \(n>0\)，\({}_{n}G\) 是在 \(S\) 上平坦且拟有限的 \(S\)-群预概形；它在 \(S\) 上有限，当且仅当 \(G\) 的 Abel 秩或约化秩在 \(S\) 上局部常值。}
\item[\textit{iii)}] \textit{对每个在 \(S\) 上可逆的整数 \(q>0\)，子群族 \({}_{q^n}G\) 在 \(G\) 中相对于 \(S\) 泛概形稠密（EGA IV，11.10.8）。}
\end{enumerate}

\SGARefLink{sga3:xvi:lemma:1:6}{引理 1.6} 适用于 \(G\)；由于 \(G\) 局部具有极大环面，\(G\) 的纤维的约化秩局部常值。因此由 \SGARefLink{sga3:xvi:lemma:1:6}{1.6}~ii)，对每个整数 \(n>0\)，\({}_{n}G\) 在 \(S\) 上有限。由\SGARefLink{sga3:viii:proposition:7:9}{第 VIII 讲，7.9}，\(u\) 为浸入。

\medskip
\noindent\textit{\SGARefLink{sga3:xvi:lemma:1:6}{1.6} 的证明。}
先考察 \(S\) 为域的谱的情形。

\medskip
\noindent\SGARefTarget{sga3:xvi:lemma:1:7}%
\textbf{引理 1.7。}---\textit{设 \(k\) 为特征 \(p\) 的域，\(G\) 为光滑、连通、幂零且幺幂秩为零的代数 \(k\)-群。则 \(G\) 是交换群，是 Abel 簇 \(A\) 以环面 \(T\) 为核的扩张。对每个整数 \(n>0\)，\({}_{n}G\) 为有限群，由秩为 \(n^{\rho_r+2\rho_{ab}}\) 的 \(k\)-代数定义；若 \((n,p)=1\)，它还是平展的。若 \((q,p)=1\)，则子群族 \({}_{q^n}G\)，\(n\in\mathbb N\)，在 \(G\) 中概形稠密。}

\smallskip
\noindent\textit{\SGARefLink{sga3:xvi:lemma:1:7}{1.7} 的证明。}
设 \(Z\) 为 \(G\) 的中心。群 \(G/Z\) 仿射（\SGARefLink{sga3:xii:6.1}{第 XII 讲，6.1}）、光滑、连通且幺幂秩为零，故为环面（BIBLE 4，定理 4）。由此 \(G\) 交换（第 XII 讲，6.4）。设 \(T\) 为 \(G\) 的唯一极大环面（第 XV 讲，3.4），则 Chevalley 定理（Séminaire Bourbaki 1956/57，第 145 号）立即表明，\(G\) 是 Abel 簇 \(A\) 以 \(T\) 为核的扩张。对每个整数 \(n>0\)，在 \(T\) 上取 \(n\) 次幂为满态射，因而有正合列
\[
  0\longrightarrow {}_{n}T\longrightarrow {}_{n}G
  \longrightarrow {}_{n}A\longrightarrow 0.
\]
Weil 的经典定理（A.~Weil，\emph{Variétés abéliennes et courbes algébriques}，第 IX 节，定理 33，推论 1）说明，\({}_{n}A\) 是由秩为 \(n^{2\rho_{ab}}\) 的 \(k\)-代数定义的有限群。由于 \({}_{n}T\) 是秩为 \(n^{\rho_r}\) 的有限群，便得到所断言的 \({}_{n}G\) 的结构。

现设 \(H\) 为对每个 \(n\) 都包含 \({}_{q^n}G\) 的 \(G\) 的最小闭子概形。由上述讨论及\SGARefLink{sga3:xv:lemma:4:6}{第 XV 讲，4.6}，\(H\) 为光滑连通子群，因而与 \(G\) 属于同一类型。在 \(H\) 上取 \(q\) 次幂为满态射（因为 \({}_{q}H\) 有限），故有正合列
\[
  0\longrightarrow {}_{q}H\longrightarrow {}_{q}G
  \longrightarrow {}_{q}(G/H)\longrightarrow 0.
\]
于是 \({}_{q}(G/H)=0\)，从而 \(G/H=0\)。因此子群 \({}_{q^n}G\) 在 \(G\) 中概形稠密。

\smallskip
\noindent\textit{续证 \SGARefLink{sga3:xvi:lemma:1:6}{1.6}~i)。}
为证明 \(G\) 交换，通常的做法先化归到 \(S\) 为 noether 概形，再化归到 \(S\) 为闭点是 \(s\) 的局部环的谱的情形。\(G\) 的中心由 \(G\) 的闭群子概形 \(Z\) 表示（\SGARefLink{sga3:xi:corollary:6:11}{第 XI 讲，6.11}）。为证明 \(Z=G\)，只需在模掉 \(S\) 的环的极大理想的每个幂后证明此式：此时 \(Z\) 将是 \(G\) 的开子群，且因 \(G\) 的纤维连通而等于 \(G\)。这将问题化归到 \(S\) 为局部 Artin 概形的情形。

设 \(q\) 为在 \(S\) 上可逆的整数。对每个整数 \(n\)，\({}_{q^n}G_s\) 是 \(G_s\) 的平展、中心、乘法型子群（\SGARefLink{sga3:xvi:lemma:1:7}{1.7}）。由于 \(G\) 光滑，\({}_{q^n}(G_s)\) 可提升为 \(G\) 的平展中心 \(S\)-群子概形 \(M_n\)（\SGARefLink{sga3:xv:proposition:1:2}{第 XV 讲，1.2}）。此时平坦子群族 \(M_n\)，\(n\in\mathbb N\)，在 \(G\) 中概形稠密（1.7 及 EGA IV，11.10.9）。由于 \(Z\) 在 \(G\) 中闭，且包含每个 \(M_n\)，便有 \(Z=G\)。

\smallskip
\noindent\textit{\SGARefLink{sga3:xvi:lemma:1:6}{1.6}~ii) 的证明。}
为证明 \({}_{n}G\) 在 \(S\) 上平坦且拟有限，只需证明 \(G\) 上的 \(n\) 次幂映射平坦且拟有限。由于 \(G\) 在 \(S\) 上平坦，可化归到 \(S\) 为域的谱的情形（EGA IV，11.3.10），而在 1.7 的证明中已指出 \(n\) 次幂映射为满态射。此外，\({}_{n}G\) 在 \(S\) 上分离，因为 \(G\) 在 \(S\) 上分离（第 VI\textsubscript{B} 讲，\SGARefLink{prop:sga3-VIB-5.2}{5.2}）。\({}_{n}G\) 在 \(S\) 上有限的充要条件是：\({}_{n}G\) 的纤维为秩局部常值的有限代数的谱。化归到 \(S\) 为 hensel 局部环的谱的情形，便可立即看出这一点。由 \SGARefLink{sga3:xvi:lemma:1:7}{1.7}，此条件等价于说 \(G\) 的纤维的 Abel 秩或约化秩局部常值。

\smallskip
\noindent\textit{\SGARefLink{sga3:xvi:lemma:1:6}{1.6}~iii) 的证明。}
为证明子群族 \({}_{q^n}G\)，\(n\in\mathbb N\)，在 \(G\) 中泛概形稠密，再次化归到 noether 情形。根据 1.7 和 1.6~ii)，只需应用 EGA IV，11.10.9。
% 底本源行 192--443 译毕。

\medskip
\noindent\textit{\SGARefLink{sga3:xvi:corollary:1:4}{1.4} b) 的证明。}
已经化归到 \(S\) 为离散赋值环的谱且 \(G^0\) 具有 Cartan 子群 \(C\) 的情形。令
\[
  N=\Norm_G(C)=\Norm_G(T),
\]
其中 \(T\) 为 \(C\) 的唯一极大环面（\SGARefLink{sga3:xii:7.1}{第 XII 讲，7.1}~a) 及 b)）。由于 \(H\) 在 \(S\) 上分离，由 \SGARefLink{sga3:xvi:lemma:1:6}{1.6}~ii) 及\SGARefLink{sga3:viii:proposition:7:12}{第 VIII 讲，7.12}，\(u\) 在 \(C\) 上的限制为闭浸入。另一方面，为证明 \(u\) 为闭浸入，只需对 \(u|_N\) 证明此性质（\SGARefLink{sga3:xvi:theorem:1:3}{1.3}~b)）。由于假设 \(W=N/C\) 由有限群概形表示，将下述引理应用于正合列
\[
  0\longrightarrow C\longrightarrow N\longrightarrow W
  \longrightarrow 1
\]
便得结论。（回忆：固有浸入为闭浸入。）

\medskip
\noindent\SGARefTarget{sga3:xvi:lemma:1:8}%
\textbf{引理 1.8。}---\textit{设 \(S\) 为预概形，\(G\) 为 \(S\) 上有限表示的 \(S\)-群预概形，是 \(S\) 上固有且有限表示的 \(S\)-群预概形 \(G''\) 以 \(S\) 上平坦且有限表示的 \(S\)-群预概形 \(G'\) 为核的扩张：}
\[
  1\longrightarrow G'\longrightarrow G\longrightarrow G''
  \longrightarrow 1.
\]
\textit{设 \(H\) 为 \(S\) 上有限表示的 \(S\)-群预概形（若 \(S\) 局部 noether，则只需局部有限表示），\(u:G\to H\) 为 \(S\)-群态射。若 \(u\) 在 \(G'\) 上的限制固有，则 \(u\) 固有。}

\smallskip
\noindent\textit{\SGARefLink{sga3:xvi:lemma:1:8}{1.8} 的证明。}
照例先化归到 \(S\) 为 noether 概形的情形（第 VI\textsubscript{B} 讲第 10 节及第 XV 讲，6.2），然后可应用固有性的赋值判据（EGA II，7.3.8）。因此设 \(S\) 为离散赋值环 \(A\) 的谱，闭点为 \(s\)，泛点为 \(t\)。设 \(\bar x\in G(t)\) 及 \(h\in H(S)\) 满足
\[
  u_t(\bar x)=h(t).
\]
必须证明 \(\bar x\) 来自唯一的元素 \(x\in G(S)\)。只需在离散赋值环 \(A\) 的忠实平坦扩张后证明 \(x\) 的存在性与唯一性。设 \(\bar y\) 为 \(\bar x\) 在 \(G''(t)\) 中的像。由于 \(G''\) 在 \(S\) 上固有，\(\bar y\) 来自唯一的元素 \(y\in G''(S)\)。设 \(X\) 为 \(y\) 在 \(G\) 中的逆像。预概形 \(X\) 在 \(S\) 上忠实平坦，因为 \(G'\) 在 \(S\) 上忠实平坦，态射 \(G\to G''\) 也忠实平坦（第 VI\textsubscript{B} 讲第 9 节）。对离散赋值环作忠实平坦更换后，可设 \(X\) 在 \(S\) 上具有截面 \(\ell\)（EGA IV，14.5.8）。将 \(\bar x\) 换为 \(\bar x\ell_t^{-1}\)，可设 \(\bar x\in G'(t)\)。此时 \(x\) 的存在性与唯一性由 \(u|_{G'}\) 的固有性得到。

\medskip
\noindent\textit{\SGARefLink{sga3:xvi:corollary:1:5}{1.5} 的证明。}

\smallskip
\noindent\textit{a)}
将在第 XIX 讲看到，若 \(G\) 约化，则 \(G\) 在忠实平坦拓扑下局部具有极大环面，Weyl 群有限，且幺幂秩为零。因此断言 a) 由 \SGARefLink{sga3:xvi:corollary:1:4}{1.4} 得到。

\smallskip
\noindent\textit{b)}
若 \(G\) 的纤维仿射且幺幂秩为零，则 \(G\) 在 fpqc 拓扑下局部具有极大环面（\SGARefLink{sga3:xv:corollary:8-18}{第 XV 讲，8.18}）。只需应用 \SGARefLink{sga3:xvi:corollary:1:4}{1.4}，并注意到：由于 \(G\) 的纤维连通、光滑、仿射且可解，其 Weyl 群为单位群（BIBLE 6，定理 1，推论 3）。

\medskip
\noindent\textit{\SGARefLink{sga3:xvi:theorem:1:3}{1.3} a) 的证明。}
由于 \(u\) 已是单态射，为证明 \(u\) 为浸入，只需证明 \(u\) 在 \(u(G)\) 的各点固有（EGA IV，15.7.1）；为此可使用局部固有性的赋值判据（EGA IV，15.7.5）。于是化归到 \(S\) 为剩余域代数闭的完备离散赋值环的谱，闭点为 \(s\)，泛点为 \(t\) 的情形。由于 \(G\) 在 \(S\) 上平坦，标准化归（参见\SGARefLink{sga3:viii:lemma:7:1}{第 VIII 讲，7.1} 的证明）使我们可设 \(H\) 在 \(S\) 上平坦，\(u_t\) 为同构，且 \(u_s\) 为底空间的同构。此时，\(u\) 是浸入等价于 \(u\) 是同构。

由假设，\(G\) 在 fpqc 拓扑下局部具有 Cartan 子群，因而可以谈论 \(G\) 的正则点构成的开子集 \(U\)（\SGARefLink{sga3:xv:theorem:7-3}{第 XV 讲，7.3}~i) 及 ii)）。

\medskip
\noindent\SGARefTarget{sga3:xvi:lemma:1:9}%
\textbf{引理 1.9。}---\textit{在上述假设下，\(u\) 是浸入的一个充分条件是 \(u|_U\) 为浸入。}

事实上，\(u|_U\) 为浸入，意即存在 \(H\) 的开子概形 \(V\) 及 \(V\) 的闭子概形 \(F\)，使 \(u|_U\) 经过 \(F\) 分解，并诱导 \(U\) 到 \(F\) 的同构。由于 \(H_t\)，从而 \(V_t\)，不可约，必有 \(V_t=F_t\)。于是 \(F\) 包含 \(V_t\) 在 \(V\) 中的概形闭包；由于 \(H\) 在 \(S\) 上平坦，此闭包等于 \(V\)。故 \(V=F\)。由此 \(u|_U\) 为开浸入，从而 \(u\) 为开浸入（第 VI\textsubscript{B} 讲，\SGARefLink{cor:sga3-VIB-2.6}{2.6}）。

余下只需证明 \(u|_U\) 为浸入。为此应用局部固有性的赋值判据。将 \(S\) 换为在 \(S\) 上忠实平坦的离散赋值环的谱后，必须证明：若 \(h\) 为 \(H\) 在 \(S\) 上的截面，且其像 \(h(S)\) 包含于 \(u(U)\)，则 \(h\) 是 \(U\) 在 \(S\) 上某截面 \(g\) 的像。只需证明 \(h\) 包含于 \(G\) 的某 Cartan 子群 \(C\) 的像中。事实上，由假设 \(u|_C\) 为浸入，故 \(h\) 为 \(C\) 的某截面 \(g\) 的像，而该截面必为 \(U\) 的截面。

设 \(a\in U(s)\) 及 \(b\in U(t)\) 分别满足
\[
  u_s(a)=h(s),
  \qquad
  u_t(b)=h(t).
\]
由于 \(u_t\) 为同构，\(h(t)\) 是 \(H_t\) 的正则点，因而包含在 \(H_t\) 的唯一 Cartan 子群 \(D_t\) 中（\SGARefLink{rem:sga3-VIB-3.2}{第 XIII 讲，3.2}）。设 \(D\) 为 \(D_t\) 在 \(H\) 中的概形闭包。

\medskip
\noindent\SGARefTarget{sga3:xvi:lemma:1:10}%
\textbf{引理 1.10。}---
\begin{enumerate}
\item[\textit{i)}] \textit{\(D_s\) 为幂零代数群（第 VI\textsubscript{B} 讲第 8 节）。}
\item[\textit{ii)}] \textit{\(\dim D_s=\nu\)，其中 \(\nu\) 为 \(G\) 的幂零秩，也是 \((H_s)_{\mathrm{red}}\) 的幂零秩。}
\item[\textit{iii)}] \textit{\(h(s)\in D_s(\kappa(s))\)。}
\end{enumerate}

\noindent\textit{i)}
由于 \(H\) 分离（第 VI\textsubscript{B} 讲，\SGARefLink{prop:sga3-VIB-5.2}{5.2}），\(D\) 也分离。此外，\(D\) 在 \(S\) 上平坦，且其泛纤维幂零。由第 VI\textsubscript{B} 讲，\SGARefLink{prop:sga3-VIB-8.4}{8.4}，\(D_s\) 幂零。

\smallskip
\noindent\textit{ii)}
由于 \((H_s)_{\mathrm{red}}\simeq G_s\)，这两个群的幂零秩相同。由于 \(G\) 在 fpqc 拓扑下局部具有 Cartan 子群，\(G_s\) 与 \(G_t\) 的幂零秩相同，记为 \(\nu\)。最后，由于 \(D\) 在 \(S\) 上平坦，
\[
  \dim D_s=\dim D_t=\nu
\]
（第 VI\textsubscript{B} 讲第 4 节）。

\smallskip
\noindent\textit{iii)}
由于 \(h_t\) 经过 \(D_t\) 分解，截面 \(h\) 显然经过 \(D\) 分解，这就证明了 iii)。

下面的引理表明，\((D_s)_{\mathrm{red}}\) 是
\[
  (H_s)_{\mathrm{red}}\simeq G_s
\]
的 Cartan 子群。

\medskip
\noindent\SGARefTarget{sga3:xvi:lemma:1:11}%
\textbf{引理 1.11。}---\textit{设 \(G\) 为代数闭域 \(k\) 上的光滑连通代数群，\(D\) 为包含正则元素 \(a\in G(k)\) 的光滑幂零代数子群。则 \(D^0\) 包含于 \(G\) 的某个 Cartan 子群中。若另有 \(\dim D\) 等于 \(G\) 的幂零秩，则 \(D\) 为 \(G\) 的 Cartan 子群，特别地，它连通。}

设 \(Z\) 为 \(G\) 的中心，令 \(G'=G/Z\)，它是仿射的（\SGARefLink{sga3:xii:6.1}{第 XII 讲，6.1}），并设 \(a'\) 及 \(D'\) 为 \(a\) 及 \(D\) 在 \(G'\) 中的像。\(G\) 的 Cartan 子群与 \(G'\) 的 Cartan 子群之间的对应（\SGARefLink{sga3:xii:6.6}{第 XII 讲，6.6}~e)）立即表明：\(a'\) 为 \(G'\) 的正则元素，且只需对 \((D',G')\) 证明本引理。因此可设 \(G\) 仿射。

设 \(s\) 为 \(a\) 的半单分量。它属于 \(D(k)\)（BIBLE 4，定理 3），且是正则的（BIBLE 7，定理 2，推论 1）。由 BIBLE 6，定理 2，\(s\) 中心化 \(D^0\)；故 \(D^0\) 包含于 \(s\) 的连通中心化子中，后者为 \(G\) 的 Cartan 子群（BIBLE 7，定理 2）。现若 \(\dim D\) 等于 \(G\) 的幂零秩，则 \(D^0\) 是 \(G\) 的 Cartan 子群，等于 \(\Centr_G(T)\)，其中 \(T\) 为 \(D^0\) 的唯一极大环面（\SGARefLink{sga3:xii:6.6}{第 XII 讲，6.6}）。但 \(D\) 幂零，因而中心化 \(T\)（BIBLE 6，定理 2），故 \(D=D^0\)。

现回到群 \(G\)。群
\[
  C_t=u_t^{-1}(D_t)
\]
为 \(G_t\) 中包含 \(b\) 的唯一 Cartan 子群。设 \(C\) 为 \(C_t\) 在 \(G\) 中的概形闭包。

由概形闭包的函子性，\(u|_C\) 经过 \(D\) 分解，因而 \(u_s(C_s)\subset D_s\)。由于 \(u_s\) 为 \(G_s\) 到 \((H_s)_{\mathrm{red}}\) 的同构，由 \SGARefLink{sga3:xvi:lemma:1:10}{1.10} 有
\[
  \dim C_s=\dim C_t=\nu=\dim D_s,
\]
且由 \SGARefLink{sga3:xvi:lemma:1:11}{1.11}，\(D_s\) 连通，故 \(u_s\) 诱导 \((C_s)_{\mathrm{red}}\) 到 \((D_s)_{\mathrm{red}}\) 的同构。因此 \((C_s)_{\mathrm{red}}\) 为 \(G_s\) 的 Cartan 子群。下述引理表明，\(C\) 本身就是 \(G\) 的 Cartan 子群。

\medskip
\noindent\SGARefTarget{sga3:xvi:lemma:1:12}%
\textbf{引理 1.12。}---\textit{设 \(S\) 为离散赋值环的谱，\(G\) 为光滑有限型 \(S\)-群预概形，\(C\) 为 \(G\) 的在 \(S\) 上平坦的群子概形，使泛纤维 \(C_t\) 为 \(G_t\) 的 Cartan 子群，而几何约化闭纤维 \((C_{\bar s})_{\mathrm{red}}\) 为 \(G_{\bar s}\) 的 Cartan 子群。则 \(C\) 为 \(G\) 的 Cartan 子群。}

必须证明 \(C\) 在 \(S\) 上光滑，而只需在基的忠实平坦扩张后证明这一点。对 \(C\) 的假设蕴含 \(G\) 的纤维的幂零秩为常数；因此可以谈论 \(G\) 的正则点构成的开子集 \(U\)（\SGARefLink{sga3:xv:theorem:7-3}{第 XV 讲，7.3}）。由于 \((C_{\bar s})_{\mathrm{red}}\) 为 \(G_{\bar s}\) 的 Cartan 子群，\(C_{\bar s}\) 包含 \(G_{\bar s}\) 的一个正则点 \(a_{\bar s}\)。由于 \(C\) 在 \(S\) 上平坦，必要时对基作忠实平坦扩张（EGA IV，14.5.8），可设 \(a_{\bar s}\) 为 \(G(s)\) 的元素，并可提升为截面 \(a\in C(S)\)。由于 \(U\) 为开集且包含 \(a(s)\)，它包含 \(a\)。设 \(C'\) 为 \(G\) 中包含 \(a\) 的唯一 Cartan 子群（\SGARefLink{sga3:xv:theorem:7-3}{第 XV 讲，7.3}）。必有 \(C_t=C'_t\)，从而 \(C'\) 与 \(C\) 均等于 \(C_t\) 在 \(G\) 中的概形闭包的单位分支（注意 \(C'\) 和 \(C\) 都在 \(S\) 上平坦）。

\smallskip
\noindent\textit{\SGARefLink{sga3:xvi:theorem:1:3}{1.3} a) 的证明完毕。}
由假设，\(u\) 在 \(G\) 的 Cartan 子群 \(C\) 上的限制为浸入。显然 \(h(S)\) 包含于 \(u(C)=D\)，故 \(h\) 为某截面 \(g\in C(S)\) 的像。
\hfill\(\square\)

\medskip
\noindent\textit{\SGARefLink{sga3:xvi:theorem:1:3}{1.3} b) 的证明。}
对浸入的情形，再次使用局部固有性的赋值判据（EGA IV，15.5）；对闭浸入的情形，则使用固有性的赋值判据（EGA II，7.3.6）。这将问题化归到 \(S\) 为剩余域代数闭的完备离散赋值环的谱的情形。设 \(s\) 为 \(S\) 的闭点，\(t\) 为泛点。将 \(H\) 换为 \(u_t(G_t)\) 在 \(H\) 中的概形闭包，可设 \(H\) 在 \(S\) 上平坦，且 \(u_t\) 为同构。

设 \(G^0\) 为 \(G\) 的连通分支。由 \SGARefLink{sga3:xvi:remark:1:2}{1.2} a)，限制 \(u|_{G^0}\) 为浸入；以 \(H^0\) 表示 \(H\) 中作为 \(G^0\) 的像的群子预概形。于是
\[
  H_s^0=(H_s)_{\mathrm{red}}^0.
\]

为验证赋值判据，必须证明：在浸入的情形下，\(H\) 在 \(S\) 上满足 \(h(S)\) 包含于 \(u(G)\) 的每个截面 \(h\)（分别在闭浸入的情形下，\(H\) 在 \(S\) 上的每个截面 \(h\)）均为 \(G\) 在 \(S\) 上某截面 \(g\) 的像。

设 \(C\) 为 \(G\) 的 Cartan 子群，并设 \(D=u(C)\) 为 \(H^0\) 中作为 \(C\) 的像的 Cartan 子群。\(S\)-群
\[
  D'=\operatorname{int}(h)D
\]
的纤维连通，故其底空间包含于 \(H^0\) 的底空间；另一方面，由于它在 \(S\) 上光滑，它是约化的，因而 \(D'\) 为 \(H^0\) 的光滑群子预概形。\(D'\) 的纤维为 \(H^0\) 的纤维的 Cartan 子群；因此 \(D'\) 为 \(H^0\) 的 Cartan 子群，且
\[
  C'=u^{-1}(D')
\]
为 \(G\) 的 Cartan 子群。但 \(\operatorname{Transpstr}_G(C,C')\) 在 \(S\) 上光滑且满射（\SGARefLink{sga3:xii:7.1}{第 XII 讲，7.1}~b)）。由于 \(S\) 为剩余域代数闭的 hensel 概形，存在 \(g\in G(S)\) 使 \(\operatorname{int}(g)C=C'\)。将 \(h\) 换为 \(u(g)^{-1}h\)，可设 \(h(t)\) 正规化 \(D_t\)（并且，在浸入情形下，\(h(s)\) 为 \(C_s\) 在 \(G_s\) 中的正规化子的某元素的像）。令
\[
  N=\Norm_G(C).
\]
由假设，\(u|_N\) 为浸入（分别为闭浸入）；故 \(h\) 确实为 \(N\) 在 \(S\) 上某截面的像，证明完成。
\hfill\(\square\)
% 底本源行 444--761 译毕。

\SGARefTarget{sga3:xvi:section:2}%
\section*{2. 商的一个可表性定理}
\addcontentsline{toc}{section}{2. 商的一个可表性定理}

先“回忆”下述结果。

\medskip
\noindent\SGARefTarget{sga3:xvi:theorem:2:1}%
\textbf{定理 2.1。}---\textit{设 \(S\) 为预概形，\(X\) 与 \(Y\) 为两个 \(S\)-预概形，\(f:X\to Y\) 为 \(S\)-态射。假设处于下述两种情形之一：}
\begin{enumerate}
\renewcommand{\labelenumi}{\textit{\alph{enumi})}}
\item \textit{态射 \(f\) 局部有限表示。}
\item \textit{预概形 \(S\) 局部 noether，且 \(X\) 在 \(S\) 上局部有限型。}
\end{enumerate}
\noindent\textit{则下述条件等价：}
\begin{enumerate}
\renewcommand{\labelenumi}{\textit{\roman{enumi})}}
\item \textit{存在 \(S\)-预概形 \(X'\) 以及 \(f\) 的分解}
\[
  f:X\xrightarrow{\,f'\,}X'\xrightarrow{\,f''\,}Y,
\]
\textit{其中 \(f'\) 为忠实平坦且局部有限表示的 \(S\)-态射，\(f''\) 为单态射。}
\item \textit{第一投影}
\[
  p_1:X\times_Y X\longrightarrow X
\]
\textit{为平坦态射。}
\end{enumerate}
\noindent\textit{此外，若上述条件成立，则 \((X',f')\) 是 \(X\) 关于 \(f\) 所定义的等价关系在 fpqc 拓扑下的商。因此 i) 中的分解 \(f=f''\circ f'\) 在同构意义下唯一。}

这个精细定理的证明将在 EGA V 中给出。也可参阅 J.-P.~Murre 的报告，\textit{Séminaire Bourbaki}，1965 年 5 月，第 294 号，定理 2、推论 2，其中处理了 \(Y\) 局部 noether 且 \(X\) 在 \(Y\) 上有限型的情形。下面将看到，可以化归到该情形。

先作几点说明。

\medskip
\noindent\textit{a)} 蕴涵 i) \(\Rightarrow\) ii) 是显然的。事实上，第一投影
\[
  p'_1:X\times_{X'}X\longrightarrow X
\]
经过 \(X\times_Y X\) 分解：
\[
  p'_1:X\times_{X'}X\xrightarrow{\,u\,}
  X\times_Y X\xrightarrow{\,p_1\,}X.
\]
因为 \(f''\) 为单态射，态射 \(u\) 为同构；因为 \(f'\) 平坦，\(p'_1\) 平坦；所以 \(p_1\) 平坦。

\medskip
\noindent\textit{b)} \SGARefLink{sga3:xvi:theorem:2:1}{2.1} 的各断言在 \(Y\) 上是局部的（因而在 \(S\) 上也是局部的）；它们在 \(X\) 上同样是局部的，这一点容易由局部有限表示的平坦态射是开映射的事实得到（EGA IV，11.3.1）。

\medskip
\noindent\textit{c)} 在 \SGARefLink{sga3:xvi:theorem:2:1}{2.1} a) 的假设下，上述讨论将问题化归到 \(X\) 与 \(Y\) 仿射且 \(f\) 有限表示的情形。将 \(S\) 换为 \(Y\)，可设 \(X\) 与 \(Y\) 在 \(S\) 上有限表示。再由 EGA IV，11.2.6 化归到 \(S\) 为 noether 概形的情形。

\medskip
\noindent\textit{d)} 在 \SGARefLink{sga3:xvi:theorem:2:1}{2.1} b) 的假设下，可设 \(S\)、\(X\) 及 \(Y\) 仿射，\(S\) 为 noether 概形，且 \(X\) 在 \(S\) 上有限型。将 \(Y\) 看作 \(S\) 上有限型仿射概形 \(Y_i\) 的滤过投射极限。概形 \(X\times_{Y_i}X\) 构成 \(X\times_S X\) 的递减滤过闭子概形族，其投射极限为 \(X\times_Y X\)。由于 \(X\times_S X\) 为 noether 概形，对一切充分大的 \(i\)，有
\[
  X\times_{Y_i}X=X\times_Y X.
\]
因此，若 \(f\) 满足 \SGARefLink{sga3:xvi:theorem:2:1}{2.1} ii) 的假设，则
\[
  f_i:X\xrightarrow{\,f\,}Y\longrightarrow Y_i
\]
也满足这些假设。由于 \(f\) 在 \(X\) 上定义的等价关系与 \(f_i\) 定义的等价关系一致，显然只需对 \(f_i\) 证明 ii) \(\Rightarrow\) i)。这将问题化归到 \(Y\) 在 \(S\) 上有限型的情形。

\subsection*{对群预概形的应用}
\addcontentsline{toc}{subsection}{对群预概形的应用}

\noindent\SGARefTarget{sga3:xvi:theorem:2:2}%
\textbf{定理 2.2。}---\textit{设 \(S\) 为预概形，\(G\) 为 \(S\) 上局部有限表示的 \(S\)-群预概形，作用于 \(S\)-预概形 \(X\)。设 \(\xi\in X(S)\) 为 \(X\) 在 \(S\) 上的截面，使 \(\xi\) 在 \(G\) 中的稳定子 \(H\) 为 \(G\) 的在 \(S\) 上平坦的 \(S\)-群子预概形。若 \(X\) 在 \(S\) 上局部有限型，或 \(S\) 局部 noether，则齐性商空间 \(G/H\) 可由 \(S\) 上局部有限表示的 \(S\)-预概形表示，且 \(S\)-态射}
\[
  f:G\longrightarrow X,\qquad g\longmapsto g\cdot\xi
\]
\textit{分解为}
\begin{center}
\begin{tikzcd}[column sep=5.6em,row sep=2.6em]
  G \arrow[d,"p"'] \arrow[dr,"f"]
    & {} \\
  G/H \arrow[r,"i"]
    & X.
\end{tikzcd}
\end{center}
\textit{其中 \(p\) 为典范投影，是忠实平坦且局部有限表示的态射，\(i\) 为单态射。（为明确起见，假定 \(G\) 在 \(X\) 上左作用。）}

\medskip
\noindent\textit{证明。}
态射 \(f\) 使 \(G\) 成为 \(X\)-预概形。由 \(\xi\) 的稳定子的定义，态射
\[
  G\times_S H\longrightarrow G\times_X G,\qquad
  (g,h)\longmapsto(g,gh)
\]
为同构。由于 \(H\) 在 \(S\) 上平坦，\(G\times_S H\) 在 \(G\) 上平坦；故第一投影
\[
  p_1:G\times_X G\longrightarrow G
\]
平坦。此外，若 \(X\) 在 \(S\) 上局部有限型，则 \(f\) 局部有限表示（EGA IV，1.4.3 v)）；否则，假设保证 \(S\) 局部 noether。因此只需将 \SGARefLink{sga3:xvi:theorem:2:1}{2.1} 应用于 \(f\)。

只需再证明 \(G/H\) 在 \(S\) 上局部有限表示，而这直接由\SGARefLink{sga3:v:proposition:9:1}{第 V 讲，9.1} 得到。
\hfill\(\square\)

\medskip
\noindent\SGARefTarget{sga3:xvi:corollary:2:3}%
\textbf{推论 2.3。}---\textit{设 \(S\) 为预概形，\(G\) 与 \(H\) 为两个 \(S\)-群预概形，\(u:G\to H\) 为 \(S\)-同态。假设 \(G\) 在 \(S\) 上局部有限表示，且 \(H\) 在 \(S\) 上局部有限型或 \(S\) 局部 noether。若 \(K=\Ker(u)\) 在 \(S\) 上平坦，则商群 \(G/K\) 可由 \(S\) 上局部有限表示的 \(S\)-群预概形表示，且 \(u\) 分解为}
\begin{center}
\begin{tikzcd}[column sep=5.8em,row sep=4.0em]
  G \arrow[rr,"u"] \arrow[dr,"p"']
    & {}
    & H \\
  {}
    & G/K \arrow[ur,"i"']
    & {}
\end{tikzcd}
\end{center}
\textit{其中 \(p\) 为典范投影，\(i\) 为单态射。}

\medskip
\noindent\textit{证明。}
取 \(X=H\)，并取 \(\xi\) 为 \(H\) 的单位截面，应用 \SGARefLink{sga3:xvi:theorem:2:2}{2.2} 即可。
\hfill\(\square\)

\medskip
\noindent\SGARefTarget{sga3:xvi:corollary:2:4}%
\textbf{推论 2.4。}---\textit{设 \(S\) 为预概形，\(G\) 为 \(S\) 上有限表示的 \(S\)-群预概形，\(H\) 为在 \(S\) 上光滑且纤维连通的 \(S\)-群预概形（因而由第 VI\textsubscript{B} 讲，\SGARefLink{cor:sga3-VIB-5.5}{5.5}，它在 \(S\) 上有限表示），\(i:H\to G\) 为 \(S\)-群单态射，从而 \(H\) 是 \(G\) 的子群。假设}
\[
  N=\Norm_G(H),
\]
\textit{即 \(G\) 的一个在 \(S\) 上有限表示的闭群子预概形所表示的正规化子（\SGARefLink{sga3:xi:corollary:6:11}{第 XI 讲，6.11}），在 \(S\) 上平坦。则 \(G/N\) 可由 \(S\) 上有限表示且在 \(S\) 上拟射影的 \(S\)-预概形表示。}

除最后一项外，所有待证断言在 \(S\) 上都是局部的。因此为证明这些断言，可设 \(S\) 拟紧，且 \(H\) 在 \(S\) 上的相对维数为常数 \(r\)。按第 XV 讲第 5 节的方法进行。对每个整数 \(n>0\)，设 \(G^{(n)}\)（分别为 \(H^{(n)}\)）为 \(G\)（分别为 \(H\)）的单位截面的第 \(n\) 个法不变量（EGA IV，16）。\(\mathcal O_S\)-模层 \(H^{(n)}\) 是 \(G^{(n)}\) 的商。由于 \(H\) 在 \(S\) 上光滑且维数为 \(r\)，\(H^{(n)}\) 典范地定义了
\[
  X_n=\operatorname{Grass}_{\varphi(n,r)}G^{(n)}
\]
的一个元素 \(\xi_n\)（EGA I，第二版，第 9 节），其中 \(\varphi(n,r)\) 为适当的整数。另一方面，\(G\) 通过表示
\[
  G\longrightarrow\operatorname{Aut}_{S\text{-gr}}(G),
  \qquad g\longmapsto\operatorname{int}(g)
\]
自然地作用于 \(G^{(n)}\)，因而也作用于 \(X_n\)。由于 \(S\) 拟紧，存在整数 \(m\)，使当 \(n\geq m\) 时，
\[
  N=\Norm_G(H)=\Norm_G H^{(n)}
  \qquad\text{（第 XI 讲，6.11 b)）。}
\]
换言之，对充分大的 \(n\)，\(N\) 为 \(\xi_n\) 的稳定子。因此 \(G/N\) 的可表性由 \SGARefLink{sga3:xvi:theorem:2:2}{2.2} 得到。\(G/N\) 在 \(S\) 上有限表示这一事实是\SGARefLink{sga3:v:proposition:9:1}{第 V 讲，9.1} 的推论。余下只需证明 \(G/N\) 在 \(S\) 上拟射影。为此，将在 \(G/N\) 上给出一个相对于 \(S\) 丰沛的典范可逆层。考虑函子
\[
  \mathscr F:(\mathbf{Sch}/S)^\circ\longrightarrow\mathbf{Ens}
\]
使对每个 \(S\)-预概形 \(T\)，有
\[
  \mathscr F(T)=
  \left\{
  \begin{array}{c}
  \text{\(T\)-子群 \(H\)，包含于 \(G_T\) 中，可表且在 \(T\) 上光滑，}\\
  \text{纤维连通}
  \end{array}
  \right\}.
\]

这样的群 \(H\) 在 \(S\) 上有限表示（第 VI\textsubscript{B} 讲，5.5），且 \(H\to G_T\) 为拟仿射态射（EGA IV，8.11.2）。由拟仿射态射的有效下降（SGA 1，VIII，7.9），\(\mathscr F\) 为 fpqc 拓扑下的层。由于 \(G/N\) 为 \(\mathscr F\) 的一个子函子的伴随层，有典范单态射
\[
  G/N\longrightarrow\mathscr F.
\]
因此存在 \(G\times_S(G/N)\) 的子群 \(H'\)，它可表，在 \(G/N\) 上光滑，纤维连通，并且对函子 \(G/N\) 是“泛”的。我断言，可逆层
\[
  \mathscr L=\bigl(\det(\operatorname{Lie}H')\bigr)^{-1}
\]
相对于 \(S\) 丰沛。按此形式，断言在 \(S\) 上是局部的，其证明与\SGARefLink{sga3:xv:remark:5:8}{第 XV 讲，5.8} 中的证明类似。

下述推论已在\SGARefLink{sga3:xiv:4.8.1}{第 XIV 讲，4.8 bis} 中预告。

\medskip
\noindent\SGARefTarget{sga3:xvi:corollary:2:5}%
\textbf{推论 2.5。}---\textit{设 \(S\) 为预概形，\(G\) 为 \(S\) 上光滑且有限表示、纤维连通的 \(S\)-群预概形，\(P\) 为 \(G\) 的抛物子群（\SGARefLink{sga3:xiv:paragraph:4-8-bis}{第 XIV 讲，4.8 bis}）。则 \(G/P\) 可由 \(S\) 上光滑且射影的 \(S\)-预概形表示。}

只需再证明（同上），\(G/P\) 可由拟射影且有限表示的 \(S\)-预概形表示。这由 \SGARefLink{sga3:xvi:corollary:2:4}{2.4} 得到，因为
\[
  P=\Norm_G(P)
\]
（同上），因而它在 \(S\) 上平坦。
\hfill\(\square\)
% 底本源行 762--1071 译毕。

\SGARefTarget{sga3:xvi:section:3}%
\section*{3. 中心平坦的群}
\addcontentsline{toc}{section}{3. 中心平坦的群}

\noindent\SGARefTarget{sga3:xvi:proposition:3:1}%
\textbf{命题 3.1。}---\textit{设 \(S\) 为预概形，\(G\) 为 \(S\) 上光滑且有限表示、纤维连通的 \(S\)-群预概形。假设 \(G\) 的中心 \(Z\) 在 \(S\) 上平坦；由\SGARefLink{sga3:xi:corollary:6:11}{第 XI 讲，6.11}，此中心可表。对每个整数 \(n>0\)，以 \(G^{(n)}\) 表示等于 \(G\) 沿单位截面的第 \(n\) 个法不变量的局部自由 \(\mathcal O_S\)-模，并以 \(\rho_n\) 表示 \(G\) 在 \(\operatorname{GL}_S(G^{(n)})\) 中的自然“伴随表示”。则：}
\begin{enumerate}
\renewcommand{\labelenumi}{\textit{\alph{enumi})}}
\item \textit{商 \(G'=G/Z\) 可由光滑、在 \(S\) 上有限表示、拟仿射且纤维仿射连通的 \(S\)-群预概形表示。}
\item \textit{表示 \(\rho_n\) 经过 \(G'\) 分解：}
\begin{center}
\begin{tikzcd}[column sep=6.0em,row sep=4.0em]
  G \arrow[rr,"\rho_n"] \arrow[dr]
    & {}
    & \operatorname{GL}_S\!\left(G^{(n)}\right) \\
  {}
    & G/Z \arrow[ur,"i_n"']
    & {}
\end{tikzcd}
\end{center}
\textit{若 \(S\) 拟紧，则对一切充分大的 \(n\)，\(i_n\) 为单态射。}
\item \textit{\(G'\) 的子环面函子 \(\mathscr T_{G'}\) 可由 \(S\) 上光滑的 \(S\)-预概形表示；若 \(S\) 拟紧，此预概形为一族在 \(S\) 上仿射的预概形之和。}
\end{enumerate}

\medskip
\noindent\textit{证明。}
a) 中的各断言在 \(S\) 上是局部的，因此可设 \(S\) 拟紧。由\SGARefLink{sga3:xi:corollary:6:11}{第 XI 讲，6.11}~b)，对一切充分大的 \(n\)，有
\[
  Z=\Ker\rho_n,
\]
故此核在 \(S\) 上平坦。\(G/Z\) 的可表性及 \(i_n\) 为单态射这一事实随即由 \SGARefLink{sga3:xvi:corollary:2:3}{2.3} 得到。由于 \(G\) 在 \(S\) 上光滑且有限表示，\(G'\) 也如此（第 VI\textsubscript{B} 讲第 9 节）。群 \(\operatorname{GL}_S(G^{(n)})\) 在 \(S\) 上仿射，而有限表示的单态射拟仿射（EGA IV，8.11.2）。因此 \(G'\) 在 \(S\) 上拟仿射，且纤维仿射。由于 \(G'\) 在 \(S\) 上光滑，函子 \(\mathscr T_{G'}\) 形式光滑（\SGARefLink{sga3:xi:proposition:2:1:bis}{第 XI 讲，2.1 bis}）。考虑到\SGARefLink{sga3:xi:remark:4:6}{第 XI 讲，4.6} 及 \SGARefLink{sga3:xvi:corollary:4:3}{4.3}，\SGARefLink{sga3:xvi:proposition:3:1}{3.1}~c) 的断言将由下述引理得到。

\medskip
\noindent\SGARefTarget{sga3:xvi:lemma:3:2}%
\textbf{引理 3.2。}---\textit{设 \(S\) 为预概形，\(G\) 与 \(H\) 为两个在 \(S\) 上有限表示的 \(S\)-群预概形，\(u:G\to H\) 为 \(S\)-群单态射；设 \(\mathscr T_G\)（分别为 \(\mathscr T_H\)）为 \(G\)（分别为 \(H\)）的子环面函子（参见第 XV 讲第 8 节）。则映射 \(T\mapsto u(T)\) 定义一个单态射（\SGARefLink{sga3:xv:lemma:8:3}{第 XV 讲，8.3}~c)）}
\[
  \widetilde u:\mathscr T_G\longrightarrow\mathscr T_H,
\]
\textit{它可由有限表示的闭浸入表示。}

按通常的做法，可化归到下述问题。设 \(T\) 为 \(H\) 的子环面，令
\[
  T'=G\times_H T
\]
为它在 \(G\) 中的逆像，设 \(u_T:T'\to T\) 为 \(u\) 诱导的 \(S\)-群单态射。必须证明，\(S\)-函子
\[
  \prod_{T/S}(T'/T)
\]
可由 \(S\) 的有限表示闭子预概形表示。但 \(T\) 是环面，故在 \(S\) 上光滑且纤维连通，只需应用\SGARefLink{sga3:xi:corollary:6:10}{第 XI 讲，6.10}。
\hfill\(\square\)

\medskip
\noindent\SGARefTarget{sga3:xvi:theorem:3:3}%
\textbf{定理 3.3。}---\textit{设 \(S\) 为预概形，\(G\) 为 \(S\) 上光滑且有限表示、纤维连通的 \(S\)-群预概形，\(Z\) 为 \(G\) 的中心。假设 \(Z\) 在 \(S\) 上平坦，且 \(G\) 的幺幂秩（\SGARefLink{sga3:xv:definition:6:1:ter}{第 XV 讲，6.1 ter}）等于 \(Z\) 的幺幂秩。则：}
\begin{enumerate}
\renewcommand{\labelenumi}{\textit{\alph{enumi})}}
\item \textit{群 \(G'=G/Z\) 可表；若 \(S\) 拟紧，则对一切充分大的 \(n\)，典范态射}
\[
  i_n:G'\longrightarrow\operatorname{GL}_S(G^{(n)})
\]
\textit{（参见 \SGARefLink{sga3:xvi:proposition:3:1}{3.1}~b)）为浸入。}
\item \textit{群 \(G'\) 在 \(S\) 上拟仿射，纤维仿射；它的中心为单位群，且 \(G'\) 的幺幂秩为零。}
\item \textit{在平展拓扑下，\(G'\) 局部具有极大环面，且这些环面也是 \(G'\) 的 Cartan 子群。\(G'\) 的极大环面函子 \(\mathscr{TM}_{G'}\)（第 XV 讲第 8 节）可由 \(S\) 上光滑且仿射的 \(S\)-预概形表示。}
\item \textit{在平展拓扑下，\(G\) 局部具有 Cartan 子群，且 \(G\) 的 Cartan 子群函子 \(\mathscr C_G\) 可由 \(S\) 上光滑且仿射的 \(S\)-预概形表示。}
\end{enumerate}

\medskip
\noindent\textit{证明。}
由 \SGARefLink{sga3:xvi:proposition:3:1}{3.1}，群 \(G'\) 由一个在 \(S\) 上光滑且拟仿射、纤维仿射的 \(S\)-预概形表示。考虑到 \(G\) 的纤维的 Cartan 子群与 \(G'\) 的纤维的 Cartan 子群之间的对应（\SGARefLink{sga3:xii:6.6}{第 XII 讲，6.6}~e)），关于 \(Z\) 的幺幂秩的假设蕴含 \(G'\) 的幺幂秩为零。现在\SGARefLink{sga3:xv:corollary:8-18}{第 XV 讲，8.18} 表明，在平展拓扑下 \(G'\) 局部具有极大环面。于是由 \SGARefLink{sga3:xvi:proposition:3:1}{3.1}~b) 及 \SGARefLink{sga3:xvi:corollary:1:4}{1.4}~a)，对一切充分大的 \(n\)，\(i_n\) 为浸入。这证明了 a)。

现证 c)。由于 \(G'\) 的幺幂秩为零，\(G'\) 的每个极大环面显然也是 \(G'\) 的 Cartan 子群。\(G'\) 的极大环面函子由 \(G'\) 的子环面函子中的一个既开又闭的子预概形表示；此外，它在 \(S\) 上有限型（\SGARefLink{sga3:xv:theorem:8-15}{第 XV 讲，8.15}）。由 \SGARefLink{sga3:xvi:proposition:3:1}{3.1}~c)，此函子由 \(S\) 上光滑且仿射的 \(S\)-预概形表示。

由于 \(G'\) 在平展拓扑下局部具有 Cartan 子群，\(G\) 也如此；且 \(G\) 的 Cartan 子群函子与 \(G'\) 的 Cartan 子群函子典范同构（\SGARefLink{sga3:xv:theorem:7-3}{第 XV 讲，7.3}~iv)）。因此 c) 蕴含 d)。

余下只需证明 b)，确切地说，即证明 \(G'\) 的中心 \(Z'\) 为单位群。由于 \(Z'\) 可表（\SGARefLink{sga3:xi:corollary:6:11}{第 XI 讲，6.11}），且在 \(S\) 上有限表示，只需证明 \(Z'\) 的纤维均为单位群（第 VI\textsubscript{B} 讲，\SGARefLink{lem:sga3-VIB-2.9}{2.9}）。于是化归到 \(S\) 为代数闭域的谱的情形。中心 \(Z'\) 显然包含于 \(G'\) 的每个 Cartan 子群，故由 c)，包含于 \(G'\) 的每个极大环面；因此它是乘法型的。另一方面，将在第 XVII 讲看到，若 \(G\) 是中心为 \(Z\) 的连通代数群，则 \(G/Z\) 的中心 \(Z'\) 幺幂。在目前的情形下，\(Z'\) 既是乘法型又是幺幂的，故为单位群（参见第 XVII 讲）。
\hfill\(\square\)

\subsection*{中心平坦的群的例子}
\addcontentsline{toc}{subsection}{中心平坦的群的例子}

\noindent\SGARefTarget{sga3:xvi:proposition:3:4}%
\textbf{命题 3.4。}---\textit{设 \(S\) 为预概形，\(G\) 为 \(S\) 上有限表示、光滑且纤维连通的 \(S\)-群预概形，\(Z\) 为 \(G\) 的中心。则在下述任一情形下，\(Z\) 在 \(S\) 上平坦：}
\begin{enumerate}
\renewcommand{\labelenumi}{\textit{\alph{enumi})}}
\item \textit{\(G\) 的纤维的幺幂秩 \(\rho_u\)（\SGARefLink{sga3:xv:definition:6:1:ter}{第 XV 讲，6.1 ter}）为零。}
\item \textit{i) \(S\) 约化；\quad ii) \(Z\) 的纤维的维数为 \(S\) 上的局部常值函数；\quad iii) \(G\) 的幺幂秩等于 \(Z\) 的幺幂秩。}
\end{enumerate}

\medskip
\noindent\textit{\SGARefLink{sga3:xvi:proposition:3:4}{3.4} a) 的证明。}
将同时证明下述命题。

\medskip
\noindent\SGARefTarget{sga3:xvi:proposition:3:5}%
\textbf{命题 3.5。}---\textit{在 \SGARefLink{sga3:xvi:proposition:3:4}{3.4}~a) 的假设下，再设 \(G\) 具有极大环面 \(T\)，并设}
\[
  C=\Centr_G(T)
\]
\textit{为 \(G\) 中与 \(T\) 对应的 Cartan 子群（\SGARefLink{sga3:xii:7.1}{第 XII 讲，7.1}）。则：}
\begin{enumerate}
\renewcommand{\labelenumi}{\textit{(\roman{enumi})}}
\item \textit{\(Z\cap T\) 为 \(G\) 的乘法型子群。}
\item \textit{\(C\) 交换且等于 \(T\cdot Z\)。}
\item \textit{若商 \(Z/(Z\cap T)\) 及 \(C/T\) 可表，则它们由 Abel 预概形（即在 \(S\) 上光滑且纤维为 Abel 簇的 \(S\)-群预概形）表示，且典范单态射}
\[
  Z/(Z\cap T)\longrightarrow C/T
\]
\textit{为同构。}
\end{enumerate}

使用第 VI\textsubscript{B} 讲第 10 节及\SGARefLink{sga3:xv:proposition:6-2}{第 XV 讲，6.2}、\SGARefLink{sga3:xvi:proposition:6:3}{6.3} 和 \SGARefLink{sga3:xv:corollary:6:3:bis}{6.3 bis} 中建立的取极限的一般性质，照例化归到 \(S\) 为 noether 概形的情形（注意 \SGARefLink{sga3:xvi:proposition:3:4}{3.4} 和 \SGARefLink{sga3:xvi:proposition:3:5}{3.5} 的各断言在 \(S\) 上是局部的）。已在 3.1 中指出，对 \(G\) 的假设蕴含 \(Z\) 可表。为证明 \(Z\) 在 \(S\) 上平坦，EGA 0\textsubscript{III}，10.2.6 将问题化归到 \(S\) 为局部 Artin 概形的情形。此时 \(G\) 光滑，故在平展拓扑下局部具有极大环面（\SGARefLink{sga3:xv:corollary:8:17}{第 XV 讲，8.17}）。因此，必要时对 \(S\) 作有限平坦扩张（证明 \(Z\) 平坦时允许这样做），可设 \(G\) 具有极大环面 \(T\)。

同样地，说明为证明 \SGARefLink{sga3:xvi:proposition:3:5}{3.5}，可化归到 \(S\) 为 Artin 概形的情形。

\medskip
\noindent i) 由于 \(Z\) 在 \(G\) 中闭（\SGARefLink{sga3:xi:corollary:6:11}{第 XI 讲，6.11}），\(Z\cap T\) 为 \(T\) 的闭群子概形。由\SGARefLink{sga3:x:corollary:4-8}{第 X 讲，4.8}~b)，\(Z\cap T\) 为乘法型的，当且仅当它在 \(S\) 上平坦。与上面一样，只需在 \(S\) 为 Artin 概形时证明 \(Z\cap T\) 平坦。

\medskip
\noindent ii) 由于 \(G\) 的幺幂秩为零，\(G\) 的每个 Cartan 子群 \(C\) 均满足\SGARefLink{sga3:xvi:lemma:1:6}{引理 1.6} 的假设，故交换。等式 \(C=T\cdot Z\) 将由 iii) 得到。

\medskip
\noindent iii) 若 \(C/T\) 可表，则 \(C/T\) 在 \(S\) 上光滑（第 VI\textsubscript{B} 讲第 9 节），且其纤维为 Abel 簇（\SGARefLink{sga3:xvi:lemma:1:7}{1.7}）；因此 \(C/T\) 为 Abel 预概形。为证明典范单态射
\[
  i:Z/(Z\cap T)\longrightarrow C/T
\]
是同构，只需在 \(S\) 为局部 Artin 概形时验证。事实上，随后可依次推出：\(Z/(Z\cap T)\) 在 \(S\) 上平坦（EGA 0\textsubscript{III}，10.2.6），\(i\) 平坦（EGA IV，11.3.10），\(i\) 为开浸入（EGA IV，17.9.1），最后 \(i\) 为同构（因为 \(C/T\) 的纤维连通）。

此后设 \(S\) 为闭点是 \(s\) 的局部 Artin 概形，且 \(G\) 具有极大环面 \(T\)。令
\[
  C=\Centr_G(T).
\]
群 \(C\) 为 \(G\) 的 Cartan 子群（\SGARefLink{sga3:xii:7.1}{第 XII 讲，7.1}），并包含 \(Z\)。

代数群 \(M_s=Z_s\cap T_s\) 为 \(G_s\) 的中心乘法型子群。由于 \(T\) 在 \(S\) 上光滑，\(M_s\) 可提升为 \(T\) 的乘法型群子概形 \(M\)（\SGARefLink{sga3:ix:theorem:3:6:bis}{第 IX 讲，3.6 bis}），并且包含于 \(G\) 的中心（\SGARefLink{sga3:ix:corollary:3:9}{第 IX 讲，3.9}），从而包含于 \(Z\cap T\)。由于 \(S\) 为 Artin 概形，且 \(M\) 与 \(T\) 在 \(S\) 上平坦，商群
\[
  (Z\cap T)/M,\qquad Z'=Z/M,\qquad A=C/T
\]
可表（第 VI\textsubscript{A} 讲第 4、5 节）。依构造，\(\bigl((Z\cap T)/M\bigr)_s\) 为单位群，故 \(Z\cap T=M\)（第 VI\textsubscript{B} 讲，\SGARefLink{lem:sga3-VIB-2.9}{2.9}）；特别地，它在 \(S\) 上平坦，这证明了 \SGARefLink{sga3:xvi:proposition:3:5}{3.5}~i)。

取商给出典范单态射
\[
  i:Z'\longrightarrow A,
\]
故它为闭浸入（第 VI\textsubscript{B} 讲，\SGARefLink{cor:sga3-VIB-1.4.2}{1.4.2}）。已经指出 \(A\) 为 Abel 预概形。只需再证明 \(i\) 为同构。这将证明 \SGARefLink{sga3:xvi:proposition:3:5}{3.5}~iii)，并蕴含 \(Z'\) 在 \(S\) 上平坦，从而 \(Z\) 在 \(S\) 上平坦（因为它是 \(Z'\) 以平坦群 \(M\) 为核的扩张；第 VI\textsubscript{B} 讲第 9 节）。由于 \(Z_s\) 与 \(G_s\) 的 Abel 秩相同（\SGARefLink{sga3:xii:6.1}{第 XII 讲，6.1}），\(i_s\) 为满态射，因而为同构。设 \(q>0\) 为在 \(S\) 上可逆的整数。由 \SGARefLink{sga3:xvi:lemma:1:6}{1.6}~iii) 的稠密性定理，为证明 \(i(Z')=A\)，只需证明对每个为 \(q\) 的幂的整数 \(n\)，\(i(Z')\) 包含 \({}_nA\)。现设 \(M_s^0\) 为 \(M_s\) 的单位分支。由对偶性立即可知，在 \(M_s^0\) 中取 \(q\) 次幂为满态射。因此，若 \(m_0\) 是 \(q\) 在 \(\operatorname{card}(M_s/M_s^0)\) 的素因数分解中的指数，则 \({}_{nm_0}Z_s\) 在
\[
  A_s\simeq Z_s/M_s
\]
中的像包含 \({}_nA_s\)。因此存在 \(Z_s\) 的乘法型子群 \(M_s(n)\)，它在 \(A_s\) 中的像为 \({}_nA_s\)。如上所述，\(M_s(n)\) 可提升为 \(G\) 的中心乘法型子群 \(M(n)\)。\(M(n)\) 在 \(A\) 中的像为乘法型子群（\SGARefLink{sga3:ix:theorem:6:8}{第 IX 讲，6.8}），且必等于 \({}_nA\)，因为在约化纤维上二者相等（\SGARefLink{sga3:ix:theorem:5:1:bis}{第 IX 讲，5.1 bis}）。至此完成了 \SGARefLink{sga3:xvi:proposition:3:4}{3.4}~a) 及 \SGARefLink{sga3:xvi:proposition:3:5}{3.5} 的证明。

\medskip
\noindent\textit{\SGARefLink{sga3:xvi:proposition:3:4}{3.4} b) 的证明。}
待证断言在 \(S\) 上是局部的，故可设 \(S\) 仿射，环为 \(A\)。将 \(A\) 看作其有限生成 \(\mathbb Z\)-子代数的归纳极限，与上面一样，化归到 \(S\) 为 noether 且约化的情形。

为证明 \(Z\) 平坦，可使用平坦性的赋值判据（EGA IV，11.8.1），这将问题化归到 \(S\) 为剩余域代数闭的完备离散赋值环的谱，泛点为 \(t\)，闭点为 \(s\) 的情形。设 \(Z'\) 为 \(Z_t\) 在 \(Z\) 中的概形闭包。必须证明 \(Z'=Z\)，甚至只需证明 \(Z_s=Z'_s\)（第 VI\textsubscript{B} 讲，\SGARefLink{cor:sga3-VIB-2.6}{2.6}）。由第 VI\textsubscript{B} 讲第 4 节，有不等式
\[
  \dim Z_t=\dim Z'_t=\dim Z'_s\leq \dim Z_s.
\]
但由假设 \(\dim Z_t=\dim Z_s\)；故 \(\dim Z'_s=\dim Z_s\)，从而 \(Z'_s\) 包含 \((Z_s)^0_{\mathrm{red}}\)。因此
\[
  G'_s=G_s/Z'_s
\]
仿射（\SGARefLink{sga3:xii:6.1}{第 XII 讲，6.1}）。考虑到 \(G\) 与 \(G'\) 的 Cartan 子群之间的对应（\SGARefLink{sga3:xii:6.6}{第 XII 讲，6.6}~e)），假设 \SGARefLink{sga3:xvi:proposition:3:4}{3.4}~b)~iii) 蕴含 \(G'_s\) 的幺幂秩为零，因而它的 Cartan 子群也就是它的极大环面。\(Z_s\) 在 \(G'_s\) 中的像 \(Z''_s\) 为 \(G'_s\) 的中心代数子群，故包含于 \(G'_s\) 的每个 Cartan 子群，即每个极大环面；因此 \(Z''_s\) 为 \(G'_s\) 的乘法型子群。在这些条件下，将在第 XVII 讲第 7 节证明，存在 \(Z_s\) 的有限乘法型子群 \(M_s\)，使其在 \(Z_s/Z'_s\) 中的像为 \(Z''_s\)（在 \(Z''_s\) 平展时，我们已在 \SGARefLink{sga3:xvi:proposition:3:4}{3.4}~a) 的证明中使用此事实）。由于 \(G\) 在 \(S\) 上光滑，而 \(S\) 是完备局部环的谱，\(M_s\) 可提升为 \(G\) 的乘法型 \(S\)-子群 \(M\)（\SGARefLink{sga3:ix:theorem:3:6:bis}{第 IX 讲，3.6 bis} 及\SGARefLink{sga3:xv:proposition:1-6}{第 XV 讲，1.6}~b)），且此子群为中心子群（\SGARefLink{sga3:ix:theorem:5:6}{第 IX 讲，5.6}~a)）。于是 \(M_t\) 包含于 \(Z'_t=Z_t\)，并且由于 \(M\) 平坦，\(M\) 包含于 \(Z'\)。因此群 \(M_s\) 包含于 \(Z'_s\)，但这蕴含 \(Z''_s\) 为单位群，即 \(Z_s=Z'_s\)。这就完成了 \SGARefLink{sga3:xvi:proposition:3:4}{3.4}~b) 的证明。

\subsection*{纤维连通而中心不平坦的光滑群预概形的例子}
\addcontentsline{toc}{subsection}{纤维连通而中心不平坦的光滑群预概形的例子}

设 \(S\) 为离散赋值环 \(A\) 的谱，\(\pi\) 为 \(A\) 的素元，\(t\) 与 \(s\) 分别为 \(S\) 的泛点与闭点。设 \(G\) 为光滑仿射 \(S\)-群，其环为
\[
  B=A[T,T^{-1},U]/F,\qquad F=1-T+\pi U,
\]
合成律定义为
\[
  (t,u),(t',u')\longmapsto (tt',\pi uu'+u+u').
\]
因此泛纤维 \(G_t\) 同构于乘法群 \((\Gm)_t\)，而闭纤维 \(G_s\) 同构于加法群 \((\Ga)_s\)。函数 \(T\) 在 \(\Gamma(G,\mathcal O_G)\) 中可逆，并定义 \(S\)-群态射
\[
  \varphi:G\longrightarrow(\Gm)_S,
\]
它在泛纤维上为同构，在闭纤维上为平凡态射。\(\varphi\) 这一数据可用来构造半直积群
\[
  H=G\cdot(\Ga)_S
\]
（注意 \((\Gm)_S\) 作用于 \((\Ga)_S\)）。\(H\) 的中心 \(Z\) 在泛纤维上为单位群，在闭纤维上等于 \(H_s\)，因而在 \(S\) 上不平坦。
% 底本源行 1072--1485 译毕。

\SGARefTarget{sga3:xvi:section:4}%
\section*{4. 纤维仿射且幺幂秩为零的群}
\addcontentsline{toc}{section}{4. 纤维仿射且幺幂秩为零的群}

\noindent\SGARefTarget{sga3:xvi:theorem:4:1}%
\textbf{定理 4.1。}---\textit{设 \(S\) 为预概形，\(G\) 为在 \(S\) 上光滑、有限表示、纤维仿射且连通、幺幂秩为零的 \(S\)-群预概形（\SGARefLink{sga3:xv:definition:6:1:ter}{第 XV 讲，6.1 ter}）。则：}
\begin{enumerate}
\renewcommand{\labelenumi}{\textit{\alph{enumi})}}
\item \textit{\(G\) 的中心 \(Z\) 为乘法型群（因而是 \(G\) 的一个约化中心；\SGARefLink{sga3:xii:4.1}{第 XII 讲，4.1}）。}
\item \textit{\(G\) 满足 \SGARefLink{sga3:xvi:theorem:3:3}{3.3} 的条件 a) 至 d)。}
\item \textit{在平展拓扑下，\(G\) 局部具有极大环面，且这些环面也是 \(G\) 的 Cartan 子群。\(G\) 的极大环面函子 \(\mathscr{TM}_G\) 可由 \(S\) 上光滑且仿射的 \(S\)-预概形表示。}
\item \textit{\(G\) 在 \(S\) 上拟仿射。此外，若 \(T\) 为 \(G\) 的极大环面，}
\[
  N=\Norm_G(T)
\]
\textit{为它在 \(G\) 中的正规化子，\(W=N/T\) 为相对于 \(T\) 的 Weyl 群，则下述条件等价：}
\begin{enumerate}
\renewcommand{\labelenumii}{\textit{(\roman{enumii})}}
\item \textit{\(G\) 在 \(S\) 上仿射；}
\item \textit{\(G'=G/Z\) 在 \(S\) 上仿射；}
\item \textit{\(N\) 在 \(S\) 上仿射；}
\item \textit{\(W\) 在 \(S\) 上仿射。}
\end{enumerate}
\textit{若 \(S\) 局部 noether 且维数至多为 \(1\)，这些条件总成立。}
\end{enumerate}

\medskip
\noindent\textit{证明。}
由于 \(G\) 的纤维仿射（因而 Abel 秩为零）且幺幂秩为零，\(G\) 在平展拓扑下局部具有极大环面（\SGARefLink{sga3:xv:corollary:8-18}{第 XV 讲，8.18}），且 \(G\) 的每个极大环面显然都是 \(G\) 的 Cartan 子群。这证明了 c) 的第一部分。为证明 \(Z\) 为乘法型，由上述结果可设 \(G\) 具有极大环面 \(T\)。由于 \(T\) 也是 Cartan 子群，它包含 \(Z\)（因为由\SGARefLink{sga3:xii:7.1}{第 XII 讲，7.1}，\(T=\Centr_G(T)\)），且
\[
  Z=Z\cap T
\]
由 \SGARefLink{sga3:xvi:proposition:3:5}{3.5}~i) 为乘法型。这证明了 a)。由 a)，断言 b) 显然成立。另一方面，函子 \(\mathscr{TM}_G\) 与 \(G\) 的 Cartan 子群函子同构（\SGARefLink{sga3:xii:7.1}{第 XII 讲，7.1}），因而由 \SGARefLink{sga3:xvi:theorem:3:3}{3.3}~c)，它在 \(S\) 上光滑且仿射。这就完成了 c) 的证明。余下只需证明 d)。

\medskip
\noindent\textit{d) 的证明。}
由于 \(G'\) 拟仿射（\SGARefLink{sga3:xvi:theorem:3:3}{3.3}~b)），而 \(Z\) 为乘法型，从而在 \(S\) 上仿射，\(G\) 拟仿射（第 VI\textsubscript{B} 讲第 9 节）。

\medskip
\noindent i) \(\Longleftrightarrow\) ii)。
若 \(G\) 仿射，则由\SGARefLink{sga3:ix:proposition:2:3}{第 IX 讲，2.3}，\(G'\) 仿射。若 \(G'\) 仿射，则 \(G\) 是仿射群以仿射群为核的扩张，故仿射（第 VI\textsubscript{B} 讲第 9 节）。

\medskip
\noindent i) \(\Longleftrightarrow\) iii)。
若 \(G\) 仿射，则 \(N\) 仿射，因为它在 \(G\) 中闭（\SGARefLink{sga3:xi:corollary:6:11}{第 XI 讲，6.11}~a)）。此外，\(G/N\) 与函子 \(\mathscr{TM}_G\) 同构（极大环面的共轭性；参见\SGARefLink{sga3:xii:7.1}{第 XII 讲，7.1}~b)），因而由 c)，它在 \(S\) 上仿射。因此，若 \(N\) 在 \(S\) 上仿射，则 \(G\) 是仿射齐性空间以仿射群为纤维的“扩张”，故仿射（第 VI\textsubscript{B} 讲第 9 节）。

\medskip
\noindent iii) \(\Longleftrightarrow\) iv)。
若 \(N\) 仿射，则 \(W=N/T\) 仿射（\SGARefLink{sga3:ix:proposition:2:3}{第 IX 讲，2.3}）；反向蕴涵由第 VI\textsubscript{B} 讲第 9 节得到。

此外，\(N/T\) 可由 \(S\) 上平展、分离且有限型的 \(S\)-群预概形表示（\SGARefLink{sga3:xv:proposition:7-1}{第 XV 讲，7.1}~iv)），故在 \(S\) 上拟有限。因此 d) 的最后一个断言由下述引理得到，其中取 \(X=W\) 及 \(Y=S\)。

\medskip
\noindent\SGARefTarget{sga3:xvi:lemma:4:2}%
\textbf{引理 4.2。}---\textit{设 \(S\) 为维数 \(1\) 的局部 noether 预概形，\(X\) 与 \(Y\) 为两个在 \(S\) 上局部有限型的 \(S\)-预概形，\(u:X\to Y\) 为拟有限且分离的 \(S\)-态射。假设对 \(S\) 的每个点 \(s\)，由 \(u\) 经基变换 \(\operatorname{Spec}\kappa(s)\to S\) 诱导的态射}
\[
  u_s:X_s\longrightarrow Y_s
\]
\textit{有限。则 \(u\) 为仿射态射。}

\medskip
待证断言在 \(Y\) 上是局部的，故可设 \(Y\)（从而 \(X\)）在 \(S\) 上有限型。由 EGA IV 第 8 节，只需在 \(S\) 为局部环 \(A\) 的谱时证明 4.2。由 fpqc 下降，可设 \(A\) 完备，进而可设 \(A\) 约化（EGA II，1.6.4），再可设 \(A\) 正规（Nagata 定理，EGA 0\textsubscript{IV} 第 22 节，及 Chevalley 定理，EGA II，6.7.1）。若 \(A\) 为域，则由假设 \(u\) 有限，因而仿射。否则 \(A\) 为离散赋值环；设 \(s\) 与 \(t\) 为 \(S\) 的闭点与泛点，\(\pi\) 为 \(A\) 的素元。设 \(y\) 为 \(Y_s\) 的点。再次应用 EGA IV 第 8 节，可将 \(Y\) 换为 \(\mathcal O_{Y,y}\) 的谱 \(Y'\)，将 \(X\) 换为
\[
  X'=X\times_Y Y';
\]
甚至可设 \(\mathcal O_{Y,y}\) 完备。由于 \(u\) 拟有限且分离，EGA II，6.2.6 表明，\(X'\) 为两个概形 \(X_1\) 与 \(X_2\) 的不交并，其中 \(X_1\) 在 \(Y'\) 上有限（因而仿射），而 \(u(X_2)\) 不包含 \(Y'\) 的闭点 \(y\)。我断言，在这些条件下 \(X_2\) 不与闭纤维 \(X'_s\) 相交。事实上，由假设，
\[
  X'_s\longrightarrow Y'_s
\]
有限，故此态射在闭子集 \((X_2)_s\) 上的限制有限，且其在 \((Y')_s\) 中的像闭。由于此像不包含局部概形 \((Y')_s\) 的闭点，\((X_2)_s\) 为空。由于 \(u_t\) 有限，态射
\[
  X'_t\longrightarrow Y'_t
\]
在 \((X_2)_t=X_2\) 上的限制有限。此外 \(Y'_t=Y'_\pi\)，故开浸入 \(Y'_t\to Y'\) 仿射。因此复合态射
\[
  X_2\longrightarrow Y'_t\longrightarrow Y'
\]
仿射，从而
\[
  X'=X_1\amalg X_2\longrightarrow Y'
\]
仿射。

\medskip
\noindent\SGARefTarget{sga3:xvi:corollary:4:3}%
\textbf{推论 4.3。}---\textit{设 \(S\) 为预概形，\(G\) 为 \(S\) 上光滑且有限表示、纤维仿射、可解、连通且幺幂秩为零的 \(S\)-群预概形。则 \(G\) 在 \(S\) 上仿射。若进一步 \(G\) 的中心为单位群且 \(S\) 拟紧，则对一切充分大的 \(n\)，典范态射}
\[
  i_n:G\longrightarrow\operatorname{GL}_S(G^{(n)})
  \qquad\textit{（参见 3.1 b)）}
\]
\textit{为闭浸入。}

\medskip
为证明 \(G\) 在 \(S\) 上仿射，可设 \(G\) 具有极大环面 \(T\)（\SGARefLink{sga3:xvi:theorem:4:1}{4.1}~c)）。由于 \(G\) 的纤维可解，相对于 \(T\) 的 Weyl 群为单位群（BIBLE，6，定理 1，推论 3），条件 \SGARefLink{sga3:xvi:theorem:4:1}{4.1}~d)~iv) 得到满足。若 \(G\) 的中心为单位群，则对一切充分大的 \(n\)，\(i_n\) 为单态射（\SGARefLink{sga3:xvi:proposition:3:1}{3.1}），因而为闭浸入（\SGARefLink{sga3:xvi:corollary:1:5}{1.5}~b)）。

\SGARefTarget{sga3:xvi:section:5}%
\section*{5. 对约化群与半单群的应用}
\addcontentsline{toc}{section}{5. 对约化群与半单群的应用}

\noindent\SGARefTarget{sga3:xvi:definition:5:1}%
\textbf{定义 5.1。}---\textit{若 \(S\)-群预概形 \(G\) 在 \(S\) 上光滑且仿射，纤维约化（分别为半单），则称 \(G\) 为\emph{约化群}（分别为\emph{半单群}）。}

\medskip
\noindent\SGARefTarget{sga3:xvi:5:1:1}%
\textbf{5.1.1。}
从第 XIX 讲起将系统研究约化群。本节需要下述性质，它们将在第 XIX 讲中证明（不使用本讲的论述）。
\begin{enumerate}
\renewcommand{\labelenumi}{\textit{\alph{enumi})}}
\item \textit{设 \(S\) 为预概形，\(G\) 为在 \(S\) 上光滑且仿射、纤维连通的 \(S\)-群预概形，\(s\) 为 \(S\) 的一个点，使 \(G_s\) 约化（分别为半单）。则存在 \(s\) 的邻域 \(U\)，使 \(G|_U\) 约化（分别为半单）。}
\item \textit{若 \(G\) 约化，则 \(G\) 在平展拓扑下局部具有极大环面；若 \(T\) 为 \(G\) 的极大环面，则 Weyl 群}
\[
  W=\Norm_G(T)/T
\]
\textit{在 \(S\) 上有限。}
\item \textit{约化群的幺幂秩为零。}
\end{enumerate}

承认这些事实后，我们来加强上述断言 a)。

\medskip
\noindent\SGARefTarget{sga3:xvi:theorem:5:2}%
\textbf{定理 5.2。}---\textit{设 \(S\) 为预概形，\(G\) 为 \(S\) 上光滑且有限表示、纤维连通的 \(S\)-群预概形，\(s\) 为 \(S\) 的一个点。则：}
\begin{enumerate}
\renewcommand{\labelenumi}{\textit{(\roman{enumi})}}
\item \textit{若 \(G\) 的纤维仿射且 \(G_s\) 约化，则存在 \(s\) 的开邻域 \(U\)，使 \(G|_U\) 约化。}
\item \textit{若 \(G_s\) 半单，则存在 \(s\) 的开邻域 \(U\)，使 \(G|_U\) 半单。}
\end{enumerate}

使用\SGARefLink{sga3:xv:proposition:6-2}{第 XV 讲，6.2}~i) 及第 VI\textsubscript{B} 讲第 10 节，化归到 \(S\) 为 noether 概形的情形。

在情形 ii) 中，考虑 \(G\) 的中心 \(Z\)，它是可表的（\SGARefLink{sga3:xi:corollary:6:11}{第 XI 讲，6.11}~a)）。由于 \(G_s\) 半单，众所周知 \(Z_s\) 有限。因此（第 VI\textsubscript{B} 讲第 4 节），存在 \(s\) 的邻域 \(U\)，使 \(Z\) 在 \(U\) 上拟有限。对 \(U\) 的每个点 \(t\)，此时 \(G_t\) 仿射（\SGARefLink{sga3:xii:6.1}{第 XII 讲，6.1}）。因此，必要时缩小 \(S\)，在情形 ii) 及情形 i) 中均可设 \(G\) 的纤维仿射。

根据 a)，为证明 \SGARefLink{sga3:xvi:theorem:5:2}{5.2}，只需证明 \(G\) 在 \(S\) 上仿射。由于 \(G\) 的纤维仿射，已知 \(G\) 的子环面函子可由光滑 \(S\)-预概形表示（\SGARefLink{sga3:xv:corollary:8:11}{第 XV 讲，8.11} 及 \SGARefLink{sga3:xv:theorem:8:8}{8.9}）。因此作覆盖 \(s\) 的平展扩张后（这是允许的），可设存在 \(G\) 的子环面 \(T\)，使 \(T_s\) 为 \(G_s\) 的极大环面。但此时
\[
  C=\Centr_G(T)
\]
为 \(G\) 的光滑子群，纤维连通，包含 \(T\)，且满足 \(C_s=T_s\)（因为 \(C_s\) 是 \(G_s\) 的 Cartan 子群，而由 c)，\(G_s\) 的幺幂秩为零）。由此 \(C=T\)，故 \(T\) 是 \(G\) 的 Cartan 子群及极大环面；特别地，\(G\) 的幺幂秩为零，\SGARefLink{sga3:xvi:theorem:4:1}{4.1} 适用。

由 \SGARefLink{sga3:xvi:theorem:4:1}{4.1}~d)，只需证明相对于 \(T\) 的 Weyl 群 \(W\) 在 \(s\) 的某邻域 \(U\) 上仿射。事实上，将证明 \(W\) 甚至在 \(s\) 的某邻域上有限。由于 \(W\) 在 \(S\) 上拟有限（\SGARefLink{sga3:xv:proposition:7-1}{第 XV 讲，7.1}~iv)），\(W\) 在 \(s\) 的某邻域 \(U\) 上有限，等价于说态射
\[
  W\longrightarrow S
\]
在 \(s\) 处固有（EGA IV，15.7 及 EGA III，4.4.2）。为证明这一点，局部固有性的赋值判据（EGA IV，15.7）将问题化归到 \(S\) 为闭点是 \(s\) 的离散赋值环的谱的情形。但此时由 \SGARefLink{sga3:xvi:theorem:4:1}{4.1}~d) 的最后一个断言，\(G\) 在 \(S\) 上仿射，故可应用上面回忆的性质 a)，断定 \(G\) 约化（分别为半单）。再应用性质 b)，可知 \(W\) 确在 \(S\) 上有限，至此完成 \SGARefLink{sga3:xvi:theorem:5:2}{5.2} 的证明。
% 底本源行 1486--1786 译毕。

\SGARefTarget{sga3:xvi:section:6}%
\section*{6. 应用：将环面的若干刚性性质推广到幺幂秩为零的群}
\addcontentsline{toc}{section}{6. 应用：将环面的若干刚性性质推广到幺幂秩为零的群}

\noindent\SGARefTarget{sga3:xvi:proposition:6:1}%
\textbf{命题 6.1。}---\textit{设 \(S\) 为预概形，\(G\) 为在 \(S\) 上光滑、有限表示、纤维连通仿射且中心为乘法型的 \(S\)-群预概形（例如 \(G\) 的幺幂秩为零；参见 \SGARefLink{sga3:xvi:theorem:4:1}{4.1}~a)）。则 \(G\) 的每个在 \(S\) 上有限表示且在 \(S\) 上拟有限的闭正规子群 \(K\)，都在 \(S\) 上有限。}

\medskip
事实上，对 \(S\) 上的每个几何点 \(\bar s\)，\((K_{\bar s})_{\mathrm{red}}\) 是 \(G_{\bar s}\) 的有限平展子群，且为正规子群，从而为中心子群（因为 \(G_{\bar s}\) 连通）。因此
\[
  K'=Z\cap K,
\]
其中 \(Z\) 表示 \(G\) 的中心，与 \(K\) 具有相同的底空间，只需证明 \(K'\) 在 \(S\) 上有限。而 \(K'\) 为乘法型群 \(Z\) 的闭拟有限子群，故有限；这一点是直接的（在 \(S\) 上局部地，\(K'\) 包含于某个适当整数 \(n\) 对应的 \({}_nZ\) 中，而 \({}_nZ\) 在 \(S\) 上有限）。

\medskip
\noindent\SGARefTarget{sga3:xvi:corollary:6:2}%
\textbf{推论 6.2。}---\textit{设 \(S\) 与 \(G\) 如上，\(K\) 为 \(G\) 的一个在 \(S\) 上有限表示、在 \(G\) 中正规且闭的 \(S\)-群子预概形，\(s\) 为 \(S\) 的一个点。若 \(K_s\) 有限（分别为单位群），则存在 \(s\) 的开邻域 \(U\)，使 \(K|_U\) 在 \(S\) 上有限（分别为单位群）。}

\medskip
由 \(K\) 的纤维维数的上半连续性（第 VI\textsubscript{B} 讲第 4 节），缩小 \(S\) 后，即可设 \(K\) 在 \(S\) 上拟有限，从而由 \SGARefLink{sga3:xvi:proposition:6:1}{6.1} 有限。若进一步 \(K_s\) 为单位群，则 Nakayama 引理容易推出 \(K\) 在 \(\operatorname{Spec}\mathcal O_s\) 上为单位群，因而在 \(s\) 的某邻域上也为单位群。

\medskip
\noindent\SGARefTarget{sga3:xvi:proposition:6:3}%
\textbf{命题 6.3。}---\textit{设 \(S\) 为预概形，\(G\) 为 \(S\) 上光滑且有限表示、纤维仿射连通且幺幂秩为零的 \(S\)-群预概形，\(H\) 为 \(S\)-群预概形，\(s\) 为 \(S\) 的一个点，并设}
\[
  u:G\longrightarrow H
\]
\textit{为 \(S\)-群态射，使 \(\Ker u_s\) 为中心子群。再假设 \(H\) 在 \(S\) 上有限表示，或 \(S\) 局部 noether 且 \(G\) 在 \(S\) 上局部有限型。则存在 \(s\) 的开邻域 \(U\)，使若令 \(K=\Ker u\)，则 \(K|_U\) 为中心乘法型子群。此外，}
\[
  G'=(G|_U)/(K|_U)
\]
\textit{可表，且由 \(u\) 取商诱导的态射}
\[
  u':G'\longrightarrow H|_U
\]
\textit{为浸入。}

\medskip
\noindent\textit{证明。}
设 \(Z\) 为 \(G\) 的中心，令 \(K'=K\cap Z\)。

\medskip
\noindent a) \(K'\) 在 \(s\) 的某邻域 \(U\) 上为乘法型子群。事实上，在 \(s\) 的某邻域 \(U\) 上作平展扩张后（这是允许的），可设 \(G\) 具有极大环面 \(T\)（\SGARefLink{sga3:xvi:theorem:4:1}{4.1}~c)）。但此时 \(T\cap K\) 为乘法型群（\SGARefLink{sga3:xv:lemma:8:3}{第 XV 讲，8.3}），其在 \(s\) 处的纤维由假设为中心子群；因此 \(T\cap K\) 在 \(s\) 的某邻域上为中心子群（\SGARefLink{sga3:ix:theorem:5:6}{第 IX 讲，5.6}~a)），从而与 \(K'\) 一致。

\medskip
\noindent b) 证明 \(K'=K\) 在 \(s\) 的某邻域 \(U\) 上成立。由 a)，已经可设 \(K'\) 为乘法型，因而在 \(S\) 上平坦。由于 \(K'_s=K_s\)，自然浸入
\[
  K'\longrightarrow K
\]
于是在 \(s\) 的某邻域 \(U\) 上为开浸入（第 VI\textsubscript{B} 讲，\SGARefLink{cor:sga3-VIB-2.6}{2.6}）。若 \(t\in U\)，则 \(K_t\) 在群
\[
  G'_t=G_t/Z_t
\]
中的像为有限平展群，在 \(G'_t\) 中正规（因为 \(K_t\) 在 \(G_t\) 中正规），故为中心子群，因而为单位群（\SGARefLink{sga3:xvi:theorem:3:3}{3.3}~b)）。这说明 \(K_t\) 包含于 \(Z_t\)，从而等于 \(K'_t\)，故在 \(U\) 上有 \(K=K'\)。

\medskip
\noindent c) 现在 \(G'\) 的可表性由 \SGARefLink{sga3:xvi:corollary:2:3}{2.3} 得到；将 \SGARefLink{sga3:xvi:theorem:4:1}{4.1}~c) 应用于 \(G'\)，由 \SGARefLink{sga3:xvi:theorem:1:3}{1.3}~a) 即知 \(u'\) 为浸入。

\medskip
\noindent\SGARefTarget{sga3:xvi:proposition:6:4}%
\textbf{命题 6.4。}---\textit{设 \(S\) 为预概形，\(G\) 为 \(S\) 上光滑且有限表示、纤维连通的 \(S\)-群预概形，\(s\) 为 \(S\) 的一个点，使 \(G_s\) 由其子环面生成（\SGARefLink{sga3:xii:8.2}{第 XII 讲，8.2}）（例如 \(G_s\) 仿射且幺幂秩为零）。设 \(H\) 为 \(S\)-群预概形，}
\[
  u,v:G\longrightarrow H
\]
\textit{为两个 \(S\)-同态，使 \(u_s=v_s\) 且 \(u_s\) 为中心同态。再假设 \(H\) 在 \(S\) 上有限表示，或 \(S\) 局部 noether。则存在 \(s\) 的邻域 \(U\)，使}
\[
  u|_U=v|_U.
\]

\medskip
照例化归到 \(S\) 为 noether 概形的情形（为研究“\(G_s\) 由其子环面生成”这一条件，使用第 VI\textsubscript{B} 讲，\SGARefLink{prop:sga3-VIB-7.4}{7.4}）。由于 \(G\) 的纤维连通，为证明在 \(s\) 的某邻域 \(U\) 上有 \(u=v\)，只需在模掉局部环 \(\mathcal O_{S,s}\) 的极大理想的每个幂后证明 \(u=v\)；这将问题化归到 \(S\) 为局部 Artin 概形的情形。

但此时 \(G\) 的极大子环面函子可由在 \(S\) 上光滑且有限型的 \(S\)-概形 \(X\) 表示（\SGARefLink{sga3:xv:corollary:8:17}{第 XV 讲，8.17}）。设 \(T\) 为 \(G_X\) 中对函子 \(X\) 泛的极大子环面。关于 \(G_s\) 的假设意即：由 \(\kappa(s)\)-态射
\[
  f:T_s\longrightarrow G_s
\]
（浸入 \(T_s\to G_{X_s}\) 与典范投影 \(G_{X_s}\to G_s\) 的复合）生成的 \(G_s\) 的代数子群（第 VI\textsubscript{B} 讲第 7 节）为整个 \(G_s\)。而 \(X_s\) 几何连通（因为若 \(\overline T\) 为 \(G_{\bar s}\) 的极大环面，\(N\) 为它在 \(G_{\bar s}\) 中的正规化子，则 \(X_{\bar s}\) 同构于 \(G_{\bar s}/N\)），且 \(T_s\) 经 \(f\) 的像包含 \(G_s\) 的单位截面。由第 VI\textsubscript{B} 讲，\SGARefLink{prop:sga3-VIB-7.4}{7.4}，存在整数 \(N>0\) 及 \(S\)-态射
\[
  f^N:T^N\longrightarrow G,
  \qquad
  T^N=T\times_S\cdots\times_S T
  \quad\text{（\(N\) 个因子），}
\]
其中 \(f^N\) 只依赖于 \(f\) 及 \(G\) 的乘法律，且 \((f^N)_s\) 满射。

考虑基变换 \(X\to S\)，以及 \(u_X\) 与 \(v_X\) 在 \(G_X\) 的子环面 \(T\) 上的限制。由假设，
\[
  u_{X_s}=v_{X_s},
\]
且
\[
  u_{X_s}:T_{X_s}\longrightarrow H_{X_s}
\]
为中心同态。于是由\SGARefLink{sga3:ix:theorem:5:1}{第 IX 讲，5.1}，有
\[
  u_X|_T=v_X|_T.
\]
写出 \(f^N\) 的明确定义，并利用 \(u\) 和 \(v\) 为同态这一事实，立即推出
\[
  uf^N=vf^N.
\]
但 \(\kappa(s)\)-态射 \(f_s^N\) 满射，且 \(G_s\) 光滑，因而约化；所以 \(f_s^N\) 泛平坦。由于 \(T\) 在 \(S\) 上光滑，因而平坦，\(T^N\) 中存在非空开子集 \(V\)，使
\[
  f^N|_V
\]
平坦（EGA IV，11.3.10）。\footnote{法文底本此处印作 \(T\)；\(f^N\) 的定义域为 \(T^N\)。}\(V\) 的像为 \(G\) 的开子集 \(W\)（EGA IV，11.3.1），且
\[
  f^N:V\longrightarrow W
\]
忠实平坦且有限表示，故为 fpqc 拓扑下的覆盖。等式 \(uf^N=vf^N\) 遂蕴含
\[
  u|_W=v|_W.
\]
由于 \(W\) 在 \(G\) 中概形稠密（EGA IV，11.10.10），而 \(H\) 在 \(S\) 上分离（第 VI\textsubscript{A} 讲，\SGARefLink{sga3:VIA:0.3}{0.3}），立即得到 \(u=v\)。
% 底本源行 1787--2008 译毕。
